%============================================================
%  Bd. 25 - Halbringe
%============================================================

\documentclass[main.tex]{subfiles}

\ifSubfilesClassLoaded{
  \BandSetupStandalone{B25}
}{
}

\title{Bd. 25 - Halbringe}
\author{Martin Kunze}
\date{}
\setcounter{file}{25}

\begin{document}

\title{Bd. 25 - Halbringe}
\author{Martin Kunze}
\date{}
\hypersetup{pageanchor=false}
\maketitle
\hypersetup{pageanchor=true}
\setcounter{tocdepth}{3}
\tableofcontents
\setcounter{file}{25}
\renewcommand{\FormulaBandID}{25}

\chapter{Einführung}

Dieser Band entwickelt Halbringe mit Eins aus einer additiven kommutativen
Monoidstruktur und einer multiplikativen Monoidstruktur. Die natürlichen
Zahlen werden als grundlegendes Beispiel und durch ihre universelle
Eigenschaft behandelt.

\chapter{Halbringe mit Eins}

\section{Definition des Halbrings mit Eins}

\FormulaDefDeltaK[Begriff des Halbrings mit Eins]{%
  \SemiringAlg{R}{+}{\cdot}{0}{1}%
}{SemiringDef}{%
  \DeltaRow{Mengen}{R\dsep 0\dsep 1\dsep a\dsep b\dsep c}
  \DeltaRow{\textbf{Axiome}}{}
  \DeltaPrem{\makecell[l]{Additives kommutatives\\ Monoid}}{%
    \CommutativeMonoidAlg{R}{+}{0}}
  \DeltaPrem{Multiplikatives Monoid}{%
    \MonoidAlg{R}{\cdot}{1}}
  \DeltaRow{Linksdistributivität}{%
    \begin{aligned}[t]
      &a,b,c\in R\vdash{}\\[-2pt]
      &a\cdot(b+c)=a\cdot b+a\cdot c
    \end{aligned}}
  \DeltaRow{Rechtsdistributivität}{%
    \begin{aligned}[t]
      &a,b,c\in R\vdash{}\\[-2pt]
      &(a+b)\cdot c=a\cdot c+b\cdot c
    \end{aligned}}
  \DeltaRow{Linksabsorption}{a\in R\vdash0\cdot a=0}
  \DeltaRow{Rechtsabsorption}{a\in R\vdash a\cdot0=0}
  \DeltaRow{\textbf{Neues Symbol}}{}
  \DeltaPrem{Halbring mit Eins}{%
    \SemiringAlg{R}{+}{\cdot}{0}{1}}
}

\FormulaAxiomDeltaK[Zugriff auf das additive kommutative Monoid]{%
  \SemiringAlg{R}{+}{\cdot}{0}{1}
  \vdash\CommutativeMonoidAlg{R}{+}{0}%
}{SemiringAdditiveCommutativeMonoidAxiom}{%
  \DeltaRow{Mengen}{R\dsep 0\dsep 1}
}

\FormulaAxiomDeltaK[Zugriff auf das multiplikative Monoid]{%
  \SemiringAlg{R}{+}{\cdot}{0}{1}
  \vdash\MonoidAlg{R}{\cdot}{1}%
}{SemiringMultiplicativeMonoidAxiom}{%
  \DeltaRow{Mengen}{R\dsep 0\dsep 1}
}

\FormulaAxiomDeltaK[Zugriff auf die Linksdistributivität]{%
  \SemiringAlg{R}{+}{\cdot}{0}{1}\dsep a,b,c\in R
  \vdash a\cdot(b+c)=a\cdot b+a\cdot c%
}{SemiringLeftDistributiveAxiom}{%
  \DeltaRow{Mengen}{R\dsep 0\dsep 1\dsep a\dsep b\dsep c}
}

\FormulaAxiomDeltaK[Zugriff auf die Rechtsdistributivität]{%
  \SemiringAlg{R}{+}{\cdot}{0}{1}\dsep a,b,c\in R
  \vdash (a+b)\cdot c=a\cdot c+b\cdot c%
}{SemiringRightDistributiveAxiom}{%
  \DeltaRow{Mengen}{R\dsep 0\dsep 1\dsep a\dsep b\dsep c}
}

\FormulaAxiomDeltaK[Zugriff auf die Linksabsorption]{%
  \SemiringAlg{R}{+}{\cdot}{0}{1}\dsep a\in R
  \vdash0\cdot a=0%
}{SemiringLeftZeroAxiom}{%
  \DeltaRow{Mengen}{R\dsep 0\dsep 1\dsep a}
}

\FormulaAxiomDeltaK[Zugriff auf die Rechtsabsorption]{%
  \SemiringAlg{R}{+}{\cdot}{0}{1}\dsep a\in R
  \vdash a\cdot0=0%
}{SemiringRightZeroAxiom}{%
  \DeltaRow{Mengen}{R\dsep 0\dsep 1\dsep a}
}

\section{Spezielle Definitionen}

\subsection{Endliche Halbringe mit Eins}

\FormulaDefDeltaK[Begriff des endlichen Halbrings mit Eins]{%
  \FiniteSemiringAlg{R}{+}{\cdot}{0_R}{1_R}%
}{FiniteSemiringDef}{%
  \DeltaRow{Mengen}{R\dsep 0_R\dsep 1_R}
  \DeltaRow{\textbf{Axiome}}{}
  \DeltaPrem{Halbring mit Eins}{\SemiringAlg{R}{+}{\cdot}{0_R}{1_R}}
  \DeltaPrem{Endlichkeit}{\Finite{R}}
  \DeltaRow{\textbf{Neues Symbol}}{}
  \DeltaPrem{Endlicher Halbring mit Eins}{%
    \FiniteSemiringAlg{R}{+}{\cdot}{0_R}{1_R}}
}

\FormulaAxiomDeltaK[Zugriff auf die Halbringstruktur]{%
  \FiniteSemiringAlg{R}{+}{\cdot}{0_R}{1_R}
  \vdash\SemiringAlg{R}{+}{\cdot}{0_R}{1_R}%
}{FiniteSemiringBaseAxiom}{%
  \DeltaRow{Mengen}{R\dsep 0_R\dsep 1_R}
}

\FormulaAxiomDeltaK[Zugriff auf die Endlichkeit]{%
  \FiniteSemiringAlg{R}{+}{\cdot}{0_R}{1_R}\vdash\Finite{R}%
}{FiniteSemiringFinitenessAxiom}{%
  \DeltaRow{Mengen}{R\dsep 0_R\dsep 1_R}
}

\subsection{Kommutative Halbringe mit Eins}

\FormulaDefDeltaK[Kommutativer Halbring mit Eins]{%
  \CommutativeSemiringAlg{R}{+}{\cdot}{0}{1}%
}{CommutativeUnitalSemiringDef}{%
  \DeltaRow{Mengen}{R\dsep 0\dsep 1}
  \DeltaRow{\textbf{Axiome}}{}
  \DeltaPrem{Halbring mit Eins}{%
    \SemiringAlg{R}{+}{\cdot}{0}{1}}
  \DeltaPrem{Kommutatives multiplikatives Monoid}{%
    \CommutativeMonoidAlg{R}{\cdot}{1}}
  \DeltaRow{\textbf{Neues Symbol}}{}
  \DeltaPrem{Kommutativer Halbring mit Eins}{%
    \CommutativeSemiringAlg{R}{+}{\cdot}{0}{1}}
}

\FormulaAxiomDeltaK[Zugriff auf den Halbring mit Eins]{%
  \CommutativeSemiringAlg{R}{+}{\cdot}{0}{1}
  \vdash\SemiringAlg{R}{+}{\cdot}{0}{1}%
}{CommutativeSemiringBaseAxiom}{%
  \DeltaRow{Mengen}{R\dsep 0\dsep 1}
}

\FormulaAxiomDeltaK[Zugriff auf das kommutative multiplikative Monoid]{%
  \CommutativeSemiringAlg{R}{+}{\cdot}{0}{1}
  \vdash\CommutativeMonoidAlg{R}{\cdot}{1}%
}{CommutativeSemiringMultiplicativeMonoidAxiom}{%
  \DeltaRow{Mengen}{R\dsep 0\dsep 1}
}

\section{Morphismen von Halbringen}

\FormulaDefDeltaK[Begriff des Halbringhomomorphismus]{%
  \SemiringHom{f}{R,\oplus,\odot,0_R,1_R}
    {S,\boxplus,\boxtimes,0_S,1_S}%
}{SemiringHomDef}{%
  \DeltaRow{Mengen}{R\dsep S\dsep0_R\dsep1_R\dsep0_S\dsep1_S}
  \DeltaRow{Funktionen}{f\dsep\oplus\dsep\odot\dsep\boxplus\dsep\boxtimes}
  \DeltaRow{\textbf{Axiome}}{}
  \DeltaPrem{Quellhalbring}{%
    \SemiringAlg{R}{\oplus}{\odot}{0_R}{1_R}}
  \DeltaPrem{Zielhalbring}{%
    \SemiringAlg{S}{\boxplus}{\boxtimes}{0_S}{1_S}}
  \DeltaPrem{Additiver Anteil}{%
    \MonoidHom{f}{R,\oplus,0_R}{S,\boxplus,0_S}}
  \DeltaPrem{Multiplikativer Anteil}{%
    \MonoidHom{f}{R,\odot,1_R}{S,\boxtimes,1_S}}
  \DeltaRow{\textbf{Neues Symbol}}{}
  \DeltaPrem{Symbol}{\mathsf{HringHom}}
}

\FormulaAxiomDeltaK[Additionserhaltung eines Halbringhomomorphismus]{%
  \begin{aligned}[t]
    &\SemiringHom{f}{R,\oplus,\odot,0_R,1_R}
      {S,\boxplus,\boxtimes,0_S,1_S}\dsep x,y\in R\\
    &\vdash f(x\oplus y)=f(x)\boxplus f(y)
  \end{aligned}
}{SemiringHomAdditionAxiom}{%
  \DeltaRow{Mengen}{R\dsep S\dsep x\dsep y}
  \DeltaRow{Funktionen}{f\dsep\oplus\dsep\odot\dsep\boxplus\dsep\boxtimes}
}

\FormulaAxiomDeltaK[Trägerabbildung eines Halbringhomomorphismus]{%
  \SemiringHom{f}{R,\oplus,\odot,0_R,1_R}
    {S,\boxplus,\boxtimes,0_S,1_S}
  \vdash f\colon R\to S%
}{SemiringHomFunctionAxiom}{%
  \DeltaRow{Mengen}{R\dsep S}
  \DeltaRow{Funktionen}{f\dsep\oplus\dsep\odot\dsep\boxplus\dsep\boxtimes}
}

\FormulaAxiomDeltaK[Multiplikationserhaltung eines Halbringhomomorphismus]{%
  \begin{aligned}[t]
    &\SemiringHom{f}{R,\oplus,\odot,0_R,1_R}
      {S,\boxplus,\boxtimes,0_S,1_S}\dsep x,y\in R\\
    &\vdash f(x\odot y)=f(x)\boxtimes f(y)
  \end{aligned}
}{SemiringHomMultiplicationAxiom}{%
  \DeltaRow{Mengen}{R\dsep S\dsep x\dsep y}
  \DeltaRow{Funktionen}{f\dsep\oplus\dsep\odot\dsep\boxplus\dsep\boxtimes}
}

\FormulaAxiomDeltaK[Null- und Einserhaltung eines Halbringhomomorphismus]{%
  \SemiringHom{f}{R,\oplus,\odot,0_R,1_R}
    {S,\boxplus,\boxtimes,0_S,1_S}
  \vdash f(0_R)=0_S\land f(1_R)=1_S%
}{SemiringHomConstantsAxiom}{%
  \DeltaRow{Mengen}{R\dsep S\dsep0_R\dsep1_R\dsep0_S\dsep1_S}
  \DeltaRow{Funktionen}{f\dsep\oplus\dsep\odot\dsep\boxplus\dsep\boxtimes}
}

\FormulaDefDeltaK[Begriff des Halbringisomorphismus]{%
  \SemiringIso{f}{R,\oplus,\odot,0_R,1_R}
    {S,\boxplus,\boxtimes,0_S,1_S}%
}{SemiringIsoDef}{%
  \DeltaRow{Mengen}{R\dsep S\dsep0_R\dsep1_R\dsep0_S\dsep1_S}
  \DeltaRow{Funktionen}{f\dsep\oplus\dsep\odot\dsep\boxplus\dsep\boxtimes}
  \DeltaRow{\textbf{Axiome}}{}
  \DeltaPrem{Homomorphismus}{%
    \SemiringHom{f}{R,\oplus,\odot,0_R,1_R}
      {S,\boxplus,\boxtimes,0_S,1_S}}
  \DeltaPrem{Bijektivität}{f\colon R\bij S}
  \DeltaRow{\textbf{Neues Symbol}}{}
  \DeltaPrem{Symbol}{\mathsf{HringIso}}
}

\subsection{Grundlegende Eigenschaften}

\FormulaThmDeltaK[Komposition von Halbringhomomorphismen]{%
  \begin{aligned}[t]
    &\SemiringHom{f}{R,\oplus,\odot,0_R,1_R}
      {S,\boxplus,\boxtimes,0_S,1_S}\dsep{}\\[-2pt]
    &\SemiringHom{g}{S,\boxplus,\boxtimes,0_S,1_S}
      {T,\uplus,\otimes,0_T,1_T}\\
    &\vdash
      \SemiringHom{g\circ f}{R,\oplus,\odot,0_R,1_R}
      {T,\uplus,\otimes,0_T,1_T}
  \end{aligned}
}{SemiringHomComposition}{%
  \DeltaRow{Mengen}{R\dsep S\dsep T}
  \DeltaRow{Funktionen}{f\dsep g\dsep\oplus\dsep\odot\dsep
    \boxplus\dsep\boxtimes\dsep\uplus\dsep\otimes}
}
\begin{tabproofwide}
  \proofstepwidestar[1]{%
    \SemiringHom{f}{R,\oplus,\odot,0_R,1_R}
      {S,\boxplus,\boxtimes,0_S,1_S}}{\rA}
  \proofstepwidestar[2]{%
    \SemiringHom{g}{S,\boxplus,\boxtimes,0_S,1_S}
      {T,\uplus,\otimes,0_T,1_T}}{\rA}
  \proofstepwidestar[1]{%
    \MonoidHom{f}{R,\oplus,0_R}{S,\boxplus,0_S}}{%
    \FormulaRefAuto{SemiringHomDef}[def]{1}}
  \proofstepwidestar[2]{%
    \MonoidHom{g}{S,\boxplus,0_S}{T,\uplus,0_T}}{%
    \FormulaRefAuto{SemiringHomDef}[def]{2}}
  \proofstepwidestar[1,2]{%
    \MonoidHom{g\circ f}{R,\oplus,0_R}{T,\uplus,0_T}}{%
    \FormulaRefAuto{MonoidHomComposition}[thm]{3,4}}
  \proofstepwidestar[1]{%
    \MonoidHom{f}{R,\odot,1_R}{S,\boxtimes,1_S}}{%
    \FormulaRefAuto{SemiringHomDef}[def]{1}}
  \proofstepwidestar[2]{%
    \MonoidHom{g}{S,\boxtimes,1_S}{T,\otimes,1_T}}{%
    \FormulaRefAuto{SemiringHomDef}[def]{2}}
  \proofstepwidestar[1,2]{%
    \MonoidHom{g\circ f}{R,\odot,1_R}{T,\otimes,1_T}}{%
    \FormulaRefAuto{MonoidHomComposition}[thm]{6,7}}
  \proofstepwidestar[1]{%
    \SemiringAlg{R}{\oplus}{\odot}{0_R}{1_R}}{%
    \FormulaRefAuto{SemiringHomDef}[def]{1}}
  \proofstepwidestar[2]{%
    \SemiringAlg{T}{\uplus}{\otimes}{0_T}{1_T}}{%
    \FormulaRefAuto{SemiringHomDef}[def]{2}}
  \proofstepwidestar[1,2]{%
    \SemiringHom{g\circ f}{R,\oplus,\odot,0_R,1_R}
      {T,\uplus,\otimes,0_T,1_T}}{%
    \FormulaRefAuto{SemiringHomDef}[def]{9,10,5,8}}
\end{tabproofwide}

\subsection{Kommutative Homomorphismen und Isomorphismen}

\FormulaDefDeltaK[Homomorphismen kommutativer Halbringe]{%
  \begin{aligned}[t]
    &\CommutativeSemiringHom{f}{R,\oplus,\odot,0_R,1_R}
      {S,\boxplus,\boxtimes,0_S,1_S}\\
    &\quad\coloneqq
      \CommutativeSemiringAlg{R}{\oplus}{\odot}{0_R}{1_R}\land
      \CommutativeSemiringAlg{S}{\boxplus}{\boxtimes}{0_S}{1_S}\\[-2pt]
    &\qquad\land
      \SemiringHom{f}{R,\oplus,\odot,0_R,1_R}
        {S,\boxplus,\boxtimes,0_S,1_S}
  \end{aligned}
}{CommutativeSemiringHomDef}{%
  \DeltaRow{Mengen}{R\dsep S\dsep0_R\dsep1_R\dsep0_S\dsep1_S}
  \DeltaRow{Funktionen}{f\dsep\oplus\dsep\odot\dsep\boxplus\dsep\boxtimes}
  \DeltaRow{\textbf{Neues Symbol}}{}
  \DeltaPrem{\makecell[l]{Homomorphismus\\ kommutativer Halbringe}}{%
    \mathsf{KHringHom}}
}

\FormulaAxiomDeltaK[Quellstruktur eines Homomorphismus kommutativer Halbringe]{%
  \begin{aligned}[t]
    &\CommutativeSemiringHom{f}{R,\oplus,\odot,0_R,1_R}
      {S,\boxplus,\boxtimes,0_S,1_S}\\[-2pt]
    &\vdash\CommutativeSemiringAlg{R}{\oplus}{\odot}{0_R}{1_R}
  \end{aligned}%
}{CommutativeSemiringHomSourceAxiom}{%
  \DeltaRow{Mengen}{R\dsep S\dsep0_R\dsep1_R\dsep0_S\dsep1_S}
  \DeltaRow{Funktionen}{f\dsep\oplus\dsep\odot\dsep\boxplus\dsep\boxtimes}
}

\FormulaAxiomDeltaK[Zielstruktur eines Homomorphismus kommutativer Halbringe]{%
  \begin{aligned}[t]
    &\CommutativeSemiringHom{f}{R,\oplus,\odot,0_R,1_R}
      {S,\boxplus,\boxtimes,0_S,1_S}\\[-2pt]
    &\vdash\CommutativeSemiringAlg{S}{\boxplus}{\boxtimes}{0_S}{1_S}
  \end{aligned}%
}{CommutativeSemiringHomTargetAxiom}{%
  \DeltaRow{Mengen}{R\dsep S\dsep0_R\dsep1_R\dsep0_S\dsep1_S}
  \DeltaRow{Funktionen}{f\dsep\oplus\dsep\odot\dsep\boxplus\dsep\boxtimes}
}

\FormulaAxiomDeltaK[Halbringhomomorphismus eines kommutativen Homomorphismus]{%
  \begin{aligned}[t]
    &\CommutativeSemiringHom{f}{R,\oplus,\odot,0_R,1_R}
      {S,\boxplus,\boxtimes,0_S,1_S}\\[-2pt]
    &\vdash
      \SemiringHom{f}{R,\oplus,\odot,0_R,1_R}
        {S,\boxplus,\boxtimes,0_S,1_S}
  \end{aligned}%
}{CommutativeSemiringHomBaseHomAxiom}{%
  \DeltaRow{Mengen}{R\dsep S\dsep0_R\dsep1_R\dsep0_S\dsep1_S}
  \DeltaRow{Funktionen}{f\dsep\oplus\dsep\odot\dsep\boxplus\dsep\boxtimes}
}

\FormulaDefDeltaK[Isomorphismen kommutativer Halbringe]{%
  \begin{aligned}[t]
    &\CommutativeSemiringIso{f}{R,\oplus,\odot,0_R,1_R}
      {S,\boxplus,\boxtimes,0_S,1_S}\\
    &\quad\coloneqq
      \CommutativeSemiringHom{f}{R,\oplus,\odot,0_R,1_R}
        {S,\boxplus,\boxtimes,0_S,1_S}\\[-2pt]
    &\qquad\land f\colon R\bij S
  \end{aligned}
}{CommutativeSemiringIsoDef}{%
  \DeltaRow{Mengen}{R\dsep S\dsep0_R\dsep1_R\dsep0_S\dsep1_S}
  \DeltaRow{Funktionen}{f\dsep\oplus\dsep\odot\dsep\boxplus\dsep\boxtimes}
  \DeltaRow{\textbf{Neues Symbol}}{}
  \DeltaPrem{Isomorphismus kommutativer Halbringe}{\mathsf{KHringIso}}
}

\FormulaAxiomDeltaK[Homomorphismus eines Isomorphismus kommutativer Halbringe]{%
  \begin{aligned}[t]
    &\CommutativeSemiringIso{f}{R,\oplus,\odot,0_R,1_R}
      {S,\boxplus,\boxtimes,0_S,1_S}\\[-2pt]
    &\vdash
      \CommutativeSemiringHom{f}{R,\oplus,\odot,0_R,1_R}
        {S,\boxplus,\boxtimes,0_S,1_S}
  \end{aligned}%
}{CommutativeSemiringIsoHomAxiom}{%
  \DeltaRow{Mengen}{R\dsep S\dsep0_R\dsep1_R\dsep0_S\dsep1_S}
  \DeltaRow{Funktionen}{f\dsep\oplus\dsep\odot\dsep\boxplus\dsep\boxtimes}
}

\FormulaAxiomDeltaK[Bijektivität eines Isomorphismus kommutativer Halbringe]{%
  \begin{aligned}[t]
    &\CommutativeSemiringIso{f}{R,\oplus,\odot,0_R,1_R}
      {S,\boxplus,\boxtimes,0_S,1_S}\\[-2pt]
    &\vdash f\colon R\bij S
  \end{aligned}%
}{CommutativeSemiringIsoBijectiveAxiom}{%
  \DeltaRow{Mengen}{R\dsep S\dsep0_R\dsep1_R\dsep0_S\dsep1_S}
  \DeltaRow{Funktionen}{f\dsep\oplus\dsep\odot\dsep\boxplus\dsep\boxtimes}
}

\chapter{Die natürlichen Zahlen als Halbring}

\section{Der kommutative Halbring mit Eins}

\FormulaThmDeltaK[Additive und multiplikative Monoidstruktur]{%
  \begin{aligned}
    &\CommutativeMonoidAlg{\mathbb{N}}{+}{0}\land{}\\[-2pt]
    &\MonoidAlg{\mathbb{N}}{\cdot}{1}
  \end{aligned}%
}{NaturalSemiringMonoidStructures}{%
  \DeltaRow{Mengen}{\mathbb{N}\dsep 0\dsep 1}
}
\begin{tabproofwide}
  \proofstepwidestar[]{%
    \CommutativeMonoidAlg{\mathbb{N}}{+}{0}}{%
    \FormulaRefAuto{NaturalAdditiveCommutativeMonoid}[thm]}
  \proofstepwidestar[]{%
    \CommutativeMonoidAlg{\mathbb{N}}{\cdot}{1}}{%
    \FormulaRefAuto{NaturalMultiplicativeCommutativeMonoid}[thm]}
  \proofstepwidestar[]{\MonoidAlg{\mathbb{N}}{\cdot}{1}}{%
    \FormulaRefAuto{CommutativeMonoidBaseMonoidAxiom}[ax]{2}}
  \proofstepwidestar[]{%
    \CommutativeMonoidAlg{\mathbb{N}}{+}{0}
    \land\MonoidAlg{\mathbb{N}}{\cdot}{1}}{%
    \rAI{1,3}}
\end{tabproofwide}

\FormulaThmDeltaK[Die natürlichen Zahlen bilden einen Halbring mit Eins]{%
  \SemiringAlg{\mathbb{N}}{+}{\cdot}{0}{1}%
}{NaturalUnitalSemiring}{%
  \DeltaRow{Mengen}{\mathbb{N}\dsep 0\dsep 1\dsep a\dsep b\dsep c}
}
\begin{tabproofwide}
  \proofstepwidestar[]{%
    \CommutativeMonoidAlg{\mathbb{N}}{+}{0}
    \land\MonoidAlg{\mathbb{N}}{\cdot}{1}}{%
    \FormulaRefAuto{NaturalSemiringMonoidStructures}[thm]}
  \proofstepwidestar[]{%
    \CommutativeMonoidAlg{\mathbb{N}}{+}{0}}{\rAEa{1}}
  \proofstepwidestar[]{%
    \MonoidAlg{\mathbb{N}}{\cdot}{1}}{\rAEb{1}}
  \proofstepwidestar[4]{a\in\mathbb{N}}{\rA}
  \proofstepwidestar[5]{b\in\mathbb{N}}{\rA}
  \proofstepwidestar[6]{c\in\mathbb{N}}{\rA}
  \proofstepwidestar[4,5,6]{%
    a\cdot(b+c)=a\cdot b+a\cdot c}{%
    \FormulaRefAuto{PeanoMulLeftDistributive}[thm]{4,5,6}}
  \proofstepwidestar[4,5,6]{%
    (a+b)\cdot c=a\cdot c+b\cdot c}{%
    \FormulaRefAuto{PeanoMulRightDistributive}[thm]{4,5,6}}
  \proofstepwidestar[4]{0\cdot a=0}{%
    \FormulaRefAuto{PeanoMulLeftZero}[thm]{4}}
  \proofstepwidestar[4]{a\cdot0=0}{%
    \FormulaRefAuto{PeanoMulRightZero}[thm]{4}}
  \proofstepwidestar[]{%
    \SemiringAlg{\mathbb{N}}{+}{\cdot}{0}{1}}{%
    \FormulaRefAuto{SemiringDef}[def]{2,3,7,8,9,10}}
\end{tabproofwide}

\FormulaThmDeltaK[Die natürlichen Zahlen bilden einen kommutativen Halbring mit Eins]{%
  \CommutativeSemiringAlg{\mathbb{N}}{+}{\cdot}{0}{1}%
}{NaturalCommutativeUnitalSemiring}{%
  \DeltaRow{Mengen}{\mathbb{N}\dsep 0\dsep 1}
}
\begin{tabproofwide}
  \proofstepwidestar[]{%
    \SemiringAlg{\mathbb{N}}{+}{\cdot}{0}{1}}{%
    \FormulaRefAuto{NaturalUnitalSemiring}[thm]}
  \proofstepwidestar[]{%
    \CommutativeMonoidAlg{\mathbb{N}}{\cdot}{1}}{%
    \FormulaRefAuto{NaturalMultiplicativeCommutativeMonoid}[thm]}
  \proofstepwidestar[]{%
    \CommutativeSemiringAlg{\mathbb{N}}{+}{\cdot}{0}{1}}{%
    \FormulaRefAuto{CommutativeUnitalSemiringDef}[def]{1,2}}
\end{tabproofwide}

\chapter{Die universelle Eigenschaft der natürlichen Zahlen}

\section{Die kanonische Abbildung}

\FormulaDefDeltaK[Nachfolgerabbildung eines Halbrings]{%
  \SemiringAlg{R}{\oplus}{\odot}{0_R}{1_R}\dsep x\in R
  \vdash
  \SemiringSuccessor{R,\oplus,\odot,0_R,1_R}(x)
  \coloneqq x\oplus1_R%
}{SemiringSuccessorDef}{%
  \DeltaRow{Mengen}{R\dsep0_R\dsep1_R\dsep x}
  \DeltaRow{Funktionen}{\oplus\dsep\odot}
  \DeltaRow{\textbf{Neues Symbol}}{}
  \DeltaPrem{Nachfolgerabbildung}{%
    \SemiringSuccessor{R,\oplus,\odot,0_R,1_R}}
}

\FormulaThmDeltaK[Die Halbringnachfolgerabbildung ist eine Endoabbildung]{%
  \SemiringAlg{R}{\oplus}{\odot}{0_R}{1_R}
  \vdash
  \SemiringSuccessor{R,\oplus,\odot,0_R,1_R}\colon R\to R%
}{SemiringSuccessorFunction}{%
  \DeltaRow{Mengen}{R\dsep0_R\dsep1_R\dsep x}
  \DeltaRow{Funktionen}{\oplus\dsep\odot}
}
\begin{tabproof}
  \proofstep{1}{\SemiringAlg{R}{\oplus}{\odot}{0_R}{1_R}}{\rA}
  \proofstep{1}{\CommutativeMonoidAlg{R}{\oplus}{0_R}}{%
    \FormulaRefAuto{SemiringAdditiveCommutativeMonoidAxiom}[ax]{1}}
  \proofstep{1}{\MonoidAlg{R}{\oplus}{0_R}}{%
    \FormulaRefAuto{CommutativeMonoidBaseMonoidAxiom}[ax]{2}}
  \proofstep{1}{\SemigroupAlg{R}{\oplus}}{%
    \FormulaRefAuto{MonoidBaseSemigroupAxiom}[ax]{3}}
  \proofstep{1}{\MonoidAlg{R}{\odot}{1_R}}{%
    \FormulaRefAuto{SemiringMultiplicativeMonoidAxiom}[ax]{1}}
  \proofstep{1}{1_R\in R}{%
    \FormulaRefAuto{MonoidIdentityCarrierAxiom}[ax]{5}}
  \proofstep{7}{x\in R}{\rA}
  \proofstep{1,7}{x\oplus1_R\in R}{%
    \FormulaRefAuto{SemigroupClosureAxiom}[ax]{4,7,6}}
  \proofstep{1,7}{%
    \SemiringSuccessor{R,\oplus,\odot,0_R,1_R}(x)=x\oplus1_R}{%
    \FormulaRefAuto{SemiringSuccessorDef}[def]{1,7}}
  \proofstep{1,7}{%
    \SemiringSuccessor{R,\oplus,\odot,0_R,1_R}(x)\in R}{%
    \rIE{9,8}}
  \proofstep{1}{%
    \SemiringSuccessor{R,\oplus,\odot,0_R,1_R}\colon R\to R}{%
    \FormulaRefAuto{%
      x\in A\vdash t(x)\coloneq y\dsep x\in A\vdash t(x)\in B
      \exists! F\colon A\to B\,\forall x\in A\,F(x)=t(x)}{9,10}}
\end{tabproof}

\FormulaThmDeltaK[Das additive Nullelement liegt im Träger]{%
  \SemiringAlg{R}{\oplus}{\odot}{0_R}{1_R}
  \vdash
  0_R\in R%
}{SemiringZeroCarrier}{%
  \DeltaRow{Mengen}{R\dsep0_R\dsep1_R}
  \DeltaRow{Funktionen}{\oplus\dsep\odot}
}
\begin{tabproof}
  \proofstep{1}{\SemiringAlg{R}{\oplus}{\odot}{0_R}{1_R}}{\rA}
  \proofstep{1}{\CommutativeMonoidAlg{R}{\oplus}{0_R}}{%
    \FormulaRefAuto{SemiringAdditiveCommutativeMonoidAxiom}[ax]{1}}
  \proofstep{1}{\MonoidAlg{R}{\oplus}{0_R}}{%
    \FormulaRefAuto{CommutativeMonoidBaseMonoidAxiom}[ax]{2}}
  \proofstep{1}{0_R\in R}{%
    \FormulaRefAuto{MonoidIdentityCarrierAxiom}[ax]{3}}
\end{tabproof}

\FormulaDefDeltaK[Kanonische Abbildung in einen Halbring]{%
  \begin{aligned}[t]
    &\SemiringAlg{R}{\oplus}{\odot}{0_R}{1_R}\vdash{}\\[-2pt]
    &\NaturalSemiringMap{R,\oplus,\odot,0_R,1_R}
      \coloneqq
      \RecFun{0_R}{\SemiringSuccessor{R,\oplus,\odot,0_R,1_R}}
  \end{aligned}
}{NaturalSemiringMapDef}{%
  \DeltaRow{Mengen}{R\dsep0_R\dsep1_R\dsep n}
  \DeltaRow{Funktionen}{\oplus\dsep\odot}
  \DeltaRow{Startwert}{}[%
    \FormulaRefAuto{SemiringZeroCarrier}[thm]]
  \DeltaRow{Schrittfunktion}{}[%
    \FormulaRefAuto{SemiringSuccessorFunction}[thm]]
  \DeltaRow{Rekursionssatz}{}[%
    \FormulaRefAuto{DedekindRecursionTheorem}[thm]]
  \DeltaRow{\textbf{Neues Symbol}}{}
  \DeltaPrem{Symbol}{\eta}
}

\FormulaThmDeltaK[Trägerabbildung der kanonischen Abbildung]{%
  \SemiringAlg{R}{\oplus}{\odot}{0_R}{1_R}
  \vdash
  \NaturalSemiringMap{R,\oplus,\odot,0_R,1_R}
    \colon\mathbb N\to R%
}{NaturalSemiringMapFunction}{%
  \DeltaRow{Mengen}{R\dsep0_R\dsep1_R}
  \DeltaRow{Funktionen}{\oplus\dsep\odot}
}
\begin{tabproof}
  \proofstep{1}{\SemiringAlg{R}{\oplus}{\odot}{0_R}{1_R}}{\rA}
  \proofstep{1}{0_R\in R}{%
    \FormulaRefAuto{SemiringZeroCarrier}[thm]{1}}
  \proofstep{1}{%
    \SemiringSuccessor{R,\oplus,\odot,0_R,1_R}\colon R\to R}{%
    \FormulaRefAuto{SemiringSuccessorFunction}[thm]{1}}
  \proofstep{1}{%
    \RecFun{0_R}{\SemiringSuccessor{R,\oplus,\odot,0_R,1_R}}
    \colon\mathbb N\to R}{%
    \FormulaRefAuto{%
      m_0\in M\dsep f\colon M\to M \vdash
      \RecFun{m_0}{f}\colon \mathbb{N}\to M}{2,3}}
  \proofstep{1}{%
    \NaturalSemiringMap{R,\oplus,\odot,0_R,1_R}
    =
    \RecFun{0_R}{\SemiringSuccessor{R,\oplus,\odot,0_R,1_R}}}{%
    \FormulaRefAuto{NaturalSemiringMapDef}[def]{1}}
  \proofstep{1}{%
    \NaturalSemiringMap{R,\oplus,\odot,0_R,1_R}
    \colon\mathbb N\to R}{\rIE{5,4}}
\end{tabproof}

\FormulaThmDeltaK[Nullwert der kanonischen Abbildung]{%
  \SemiringAlg{R}{\oplus}{\odot}{0_R}{1_R}
  \vdash
  \NaturalSemiringMap{R,\oplus,\odot,0_R,1_R}(0)=0_R%
}{NaturalSemiringMapZero}{%
  \DeltaRow{Mengen}{R\dsep0_R\dsep1_R}
  \DeltaRow{Funktionen}{\oplus\dsep\odot}
}
\begin{tabproof}
  \proofstep{1}{\SemiringAlg{R}{\oplus}{\odot}{0_R}{1_R}}{\rA}
  \proofstep{1}{0_R\in R}{%
    \FormulaRefAuto{SemiringZeroCarrier}[thm]{1}}
  \proofstep{1}{%
    \SemiringSuccessor{R,\oplus,\odot,0_R,1_R}\colon R\to R}{%
    \FormulaRefAuto{SemiringSuccessorFunction}[thm]{1}}
  \proofstep{1}{%
    \RecFun{0_R}{\SemiringSuccessor{R,\oplus,\odot,0_R,1_R}}(0)
    =0_R}{%
    \FormulaRefAuto{%
      m_0\in M\dsep f\colon M\to M \vdash
      \RecFun{m_0}{f}(0)=m_0}{2,3}}
  \proofstep{1}{%
    \NaturalSemiringMap{R,\oplus,\odot,0_R,1_R}
    =
    \RecFun{0_R}{\SemiringSuccessor{R,\oplus,\odot,0_R,1_R}}}{%
    \FormulaRefAuto{NaturalSemiringMapDef}[def]{1}}
  \proofstep{1}{%
    \NaturalSemiringMap{R,\oplus,\odot,0_R,1_R}(0)=0_R}{%
    \rIE{5,4}}
\end{tabproof}

\FormulaThmDeltaK[Nachfolgerwert der kanonischen Abbildung]{%
  \begin{aligned}[t]
    &\SemiringAlg{R}{\oplus}{\odot}{0_R}{1_R}\dsep n\in\mathbb N\\
    &\vdash
      \NaturalSemiringMap{R,\oplus,\odot,0_R,1_R}(\Succ(n))
      =
      \NaturalSemiringMap{R,\oplus,\odot,0_R,1_R}(n)\oplus1_R
  \end{aligned}
}{NaturalSemiringMapSuccessor}{%
  \DeltaRow{Mengen}{R\dsep0_R\dsep1_R\dsep n}
  \DeltaRow{Funktionen}{\oplus\dsep\odot}
}
\begin{notation*}
\[
  \eta_R:=\NaturalSemiringMap{R,\oplus,\odot,0_R,1_R},
  \qquad
  s_R:=\SemiringSuccessor{R,\oplus,\odot,0_R,1_R}.
\]
\end{notation*}
\begin{tabproofwide}
  \proofstepwidestar[1]{%
    \SemiringAlg{R}{\oplus}{\odot}{0_R}{1_R}}{\rA}
  \proofstepwidestar[2]{n\in\mathbb N}{\rA}
  \proofstepwidestar[1]{%
    0_R\in R}{%
    \FormulaRefAuto{SemiringZeroCarrier}[thm]{1}}
  \proofstepwidestar[1]{s_R\colon R\to R}{%
    \FormulaRefAuto{SemiringSuccessorFunction}[thm]{1}}
  \proofstepwidestar[1]{\eta_R=\RecFun{0_R}{s_R}}{%
    \FormulaRefAuto{NaturalSemiringMapDef}[def]{1}}
  \proofstepwidestar[1]{\eta_R\colon\mathbb N\to R}{%
    \FormulaRefAuto{NaturalSemiringMapFunction}[thm]{1}}
  \proofstepwidestar[1,2]{\eta_R(n)\in R}{%
    \FormulaRefAuto{F\colon A\to B\dsep a\in A\vdash F(a)\in B}{6,2}}
  \proofstepwidestar[1]{\RecFun{0_R}{s_R}=\eta_R}{%
    \FormulaRefAuto{a=b\vdash b=a}{5}}

  \proofstepwide[1,2]{\eta_R(\Succ(n))}{=}{%
    \RecFun{0_R}{s_R}(\Succ(n))}{\rIE{5}}
  \proofstepwide[1,2]{}{=}{%
    s_R\!\left(\RecFun{0_R}{s_R}(n)\right)}{%
    \FormulaRefAuto{%
      m_0\in M\dsep f\colon M\to M\dsep n\in\mathbb{N}\vdash
      \RecFun{m_0}{f}(\Succ(n))=f(\RecFun{m_0}{f}(n))}{3,4,2}}
  \proofstepwide[1,2]{}{=}{s_R\!\left(\eta_R(n)\right)}{\rIE{8}}
  \proofstepwide[1,2]{}{=}{\eta_R(n)\oplus1_R}{%
    \FormulaRefAuto{SemiringSuccessorDef}[def]{1,7}}
  \proofstepwidestar[1,2]{%
    \eta_R(\Succ(n))=\eta_R(n)\oplus1_R}{%
    \rChain{9,12}}
\end{tabproofwide}

\FormulaThmDeltaK[Einserhaltung der kanonischen Abbildung]{%
  \SemiringAlg{R}{\oplus}{\odot}{0_R}{1_R}
  \vdash
  \NaturalSemiringMap{R,\oplus,\odot,0_R,1_R}(1)=1_R%
}{NaturalSemiringMapOne}{%
  \DeltaRow{Mengen}{R\dsep0_R\dsep1_R}
  \DeltaRow{Funktionen}{\oplus\dsep\odot}
}
\begin{tabproofwide}
  \proofstepwidestar[1]{%
    \SemiringAlg{R}{\oplus}{\odot}{0_R}{1_R}}{\rA}
  \proofstepwidestar[]{0\in\mathbb N}{\FormulaRefAuto{PeanoZero}[ax]}
  \proofstepwidestar[]{1=\Succ(0)}{\FormulaRefAuto{PeanoOne}[def]}
  \proofstepwidestar[1]{%
    \NaturalSemiringMap{R,\oplus,\odot,0_R,1_R}(0)=0_R}{%
    \FormulaRefAuto{NaturalSemiringMapZero}[thm]{1}}
  \proofstepwidestar[1]{\CommutativeMonoidAlg{R}{\oplus}{0_R}}{%
    \FormulaRefAuto{SemiringAdditiveCommutativeMonoidAxiom}[ax]{1}}
  \proofstepwidestar[1]{\MonoidAlg{R}{\oplus}{0_R}}{%
    \FormulaRefAuto{CommutativeMonoidBaseMonoidAxiom}[ax]{5}}
  \proofstepwidestar[1]{\MonoidAlg{R}{\odot}{1_R}}{%
    \FormulaRefAuto{SemiringMultiplicativeMonoidAxiom}[ax]{1}}
  \proofstepwidestar[1]{1_R\in R}{%
    \FormulaRefAuto{MonoidIdentityCarrierAxiom}[ax]{7}}

  \proofstepwide[1]{%
    \NaturalSemiringMap{R,\oplus,\odot,0_R,1_R}(1)}{=}{%
    \NaturalSemiringMap{R,\oplus,\odot,0_R,1_R}(\Succ(0))}{\rIE{3}}
  \proofstepwide[1]{}{=}{%
    \NaturalSemiringMap{R,\oplus,\odot,0_R,1_R}(0)\oplus1_R}{%
    \FormulaRefAuto{NaturalSemiringMapSuccessor}[thm]{1,2}}
  \proofstepwide[1]{}{=}{0_R\oplus1_R}{\rIE{4}}
  \proofstepwide[1]{}{=}{1_R}{%
    \FormulaRefAuto{MonoidLeftIdentityAxiom}[ax]{6,8}}
  \proofstepwidestar[1]{%
    \NaturalSemiringMap{R,\oplus,\odot,0_R,1_R}(1)=1_R}{%
    \rChain{9,12}}
\end{tabproofwide}

\begin{notation*}
  \[
    \eta_R\coloneqq
    \NaturalSemiringMap{R,\oplus,\odot,0_R,1_R}.
  \]
\end{notation*}

\FormulaThmDeltaK[Nullfall der Additionserhaltung]{%
  \begin{aligned}[t]
    &\SemiringAlg{R}{\oplus}{\odot}{0_R}{1_R}\dsep m\in\mathbb N\\
    &\vdash
      \NaturalSemiringMap{R,\oplus,\odot,0_R,1_R}(m+0)
      =
      \NaturalSemiringMap{R,\oplus,\odot,0_R,1_R}(m)
      \oplus
      \NaturalSemiringMap{R,\oplus,\odot,0_R,1_R}(0)
  \end{aligned}
}{NaturalSemiringMapAdditionZero}{%
  \DeltaRow{Mengen}{R\dsep0_R\dsep1_R\dsep m}
  \DeltaRow{Funktionen}{\oplus\dsep\odot}
}
\begin{tabproofwide}
  \proofstepwidestar[1]{%
    \SemiringAlg{R}{\oplus}{\odot}{0_R}{1_R}}{\rA}
  \proofstepwidestar[2]{m\in\mathbb N}{\rA}
  \proofstepwidestar[1]{%
    \CommutativeMonoidAlg{R}{\oplus}{0_R}}{%
    \FormulaRefAuto{SemiringAdditiveCommutativeMonoidAxiom}[ax]{1}}
  \proofstepwidestar[1]{\MonoidAlg{R}{\oplus}{0_R}}{%
    \FormulaRefAuto{CommutativeMonoidBaseMonoidAxiom}[ax]{3}}
  \proofstepwidestar[1]{\eta_R\colon\mathbb N\to R}{%
    \FormulaRefAuto{NaturalSemiringMapFunction}[thm]{1}}
  \proofstepwidestar[1,2]{\eta_R(m)\in R}{%
    \FormulaRefAuto{F\colon A\to B\dsep a\in A\vdash F(a)\in B}{5,2}}
  \proofstepwide[1,2]{\eta_R(m)\oplus0_R}{=}{\eta_R(m)}{%
    \FormulaRefAuto{MonoidRightIdentityAxiom}[ax]{4,6}}
  \proofstepwidestar[1]{\eta_R(0)=0_R}{%
    \FormulaRefAuto{NaturalSemiringMapZero}[thm]{1}}
  \proofstepwidestar[1]{0_R=\eta_R(0)}{%
    \FormulaRefAuto{a=b\vdash b=a}{8}}
  \proofstepwide[2]{\eta_R(m+0)}{=}{\eta_R(m)}{%
    \FormulaRefAuto{PeanoAddRightZero}[thm]{2}}
  \proofstepwide[1,2]{}{=}{\eta_R(m)\oplus0_R}{%
    \FormulaRefAuto{a=b\vdash b=a}{7}}
  \proofstepwide[1,2]{}{=}{\eta_R(m)\oplus\eta_R(0)}{\rIE{9}}
  \proofstepwidestar[1,2]{%
    \eta_R(m+0)=\eta_R(m)\oplus\eta_R(0)}{%
    \rChain{10,12}}
\end{tabproofwide}

\FormulaThmDeltaK[Nachfolgerschritt der Additionserhaltung]{%
  \begin{aligned}[t]
    &\SemiringAlg{R}{\oplus}{\odot}{0_R}{1_R}\dsep m,n\in\mathbb N\dsep{}\\[-2pt]
    &\NaturalSemiringMap{R,\oplus,\odot,0_R,1_R}(m+n)
      =
      \NaturalSemiringMap{R,\oplus,\odot,0_R,1_R}(m)
      \oplus
      \NaturalSemiringMap{R,\oplus,\odot,0_R,1_R}(n)\\
    &\vdash
      \NaturalSemiringMap{R,\oplus,\odot,0_R,1_R}(m+\Succ(n))
      =
      \NaturalSemiringMap{R,\oplus,\odot,0_R,1_R}(m)
      \oplus
      \NaturalSemiringMap{R,\oplus,\odot,0_R,1_R}(\Succ(n))
  \end{aligned}
}{NaturalSemiringMapAdditionStep}{%
  \DeltaRow{Mengen}{R\dsep0_R\dsep1_R\dsep m\dsep n}
  \DeltaRow{Funktionen}{\oplus\dsep\odot}
}
\begin{tabproofsplitwide}
  \proofpartwideindR[Assoziativität im Zielhalbring]{%
    \begin{aligned}[t]
      &\SemiringAlg{R}{\oplus}{\odot}{0_R}{1_R}\dsep
        m,n\in\mathbb N\vdash{}\\[-2pt]
      &\bigl(\eta_R(m)\oplus\eta_R(n)\bigr)\oplus1_R
        =\eta_R(m)\oplus\bigl(\eta_R(n)\oplus1_R\bigr)
    \end{aligned}
  }{%
    \SemiringAlg{R}{\oplus}{\odot}{0_R}{1_R}\dsep
    m,n\in\mathbb N
    \vdash \bigl(\eta_R(m)\oplus\eta_R(n)\bigr)\oplus1_R
      =\eta_R(m)\oplus\bigl(\eta_R(n)\oplus1_R\bigr)%
  }[NaturalSemiringMapAdditionStepAssociativity]
    \proofstepwidestar[1]{%
      \SemiringAlg{R}{\oplus}{\odot}{0_R}{1_R}}{\rA}
    \proofstepwidestar[2]{m\in\mathbb N}{\rA}
    \proofstepwidestar[3]{n\in\mathbb N}{\rA}
    \proofstepwidestar[1]{%
      \CommutativeMonoidAlg{R}{\oplus}{0_R}}{%
      \FormulaRefAuto{SemiringAdditiveCommutativeMonoidAxiom}[ax]{1}}
    \proofstepwidestar[1]{\MonoidAlg{R}{\oplus}{0_R}}{%
      \FormulaRefAuto{CommutativeMonoidBaseMonoidAxiom}[ax]{4}}
    \proofstepwidestar[1]{\SemigroupAlg{R}{\oplus}}{%
      \FormulaRefAuto{MonoidBaseSemigroupAxiom}[ax]{5}}
    \proofstepwidestar[1]{\eta_R\colon\mathbb N\to R}{%
      \FormulaRefAuto{NaturalSemiringMapFunction}[thm]{1}}
    \proofstepwidestar[1,2]{\eta_R(m)\in R}{%
      \FormulaRefAuto{F\colon A\to B\dsep a\in A\vdash F(a)\in B}{7,2}}
    \proofstepwidestar[1,3]{\eta_R(n)\in R}{%
      \FormulaRefAuto{F\colon A\to B\dsep a\in A\vdash F(a)\in B}{7,3}}
    \proofstepwidestar[1]{\MonoidAlg{R}{\odot}{1_R}}{%
      \FormulaRefAuto{SemiringMultiplicativeMonoidAxiom}[ax]{1}}
    \proofstepwidestar[1]{1_R\in R}{%
      \FormulaRefAuto{MonoidIdentityCarrierAxiom}[ax]{10}}
    \proofstepwidestar[1,2,3]{%
      \bigl(\eta_R(m)\oplus\eta_R(n)\bigr)\oplus1_R
      =\eta_R(m)\oplus\bigl(\eta_R(n)\oplus1_R\bigr)}{%
      \FormulaRefAuto{SemigroupAssociativityAxiom}[ax]{6,8,9,11}}
  \closeproofpartwideindR

  \proofpartwideindR[Gleichheitskette des Nachfolgerschritts]{%
    \begin{aligned}[t]
      &\SemiringAlg{R}{\oplus}{\odot}{0_R}{1_R}\dsep m,n\in\mathbb N\dsep{}\\[-2pt]
      &\eta_R(m+n)=\eta_R(m)\oplus\eta_R(n)\vdash{}\\[-2pt]
      &\eta_R(m+\Succ(n))=\eta_R(m)\oplus\eta_R(\Succ(n))
    \end{aligned}
  }{%
    \SemiringAlg{R}{\oplus}{\odot}{0_R}{1_R}\dsep m,n\in\mathbb N\dsep
    \eta_R(m+n)=\eta_R(m)\oplus\eta_R(n)
    \vdash \eta_R(m+\Succ(n))=\eta_R(m)\oplus\eta_R(\Succ(n))%
  }[NaturalSemiringMapAdditionStepConclusion]
    \proofstepwidestar[1]{%
      \SemiringAlg{R}{\oplus}{\odot}{0_R}{1_R}}{\rA}
    \proofstepwidestar[2]{m\in\mathbb N}{\rA}
    \proofstepwidestar[3]{n\in\mathbb N}{\rA}
    \proofstepwidestar[4]{\eta_R(m+n)=\eta_R(m)\oplus\eta_R(n)}{\rA}
    \proofstepwidestar[2,3]{m+n\in\mathbb N}{%
      \FormulaRefAuto{PeanoAddClosure}[thm]{2,3}}
    \proofstepwidestar[1,2,3]{%
      \bigl(\eta_R(m)\oplus\eta_R(n)\bigr)\oplus1_R
      =\eta_R(m)\oplus\bigl(\eta_R(n)\oplus1_R\bigr)}{%
      \FormulaRefAuto{NaturalSemiringMapAdditionStepAssociativity}[thm]{1,2,3}}
    \proofstepwidestar[1,3]{\eta_R(\Succ(n))=\eta_R(n)\oplus1_R}{%
      \FormulaRefAuto{NaturalSemiringMapSuccessor}[thm]{1,3}}
    \proofstepwidestar[1,3]{\eta_R(n)\oplus1_R=\eta_R(\Succ(n))}{%
      \FormulaRefAuto{a=b\vdash b=a}{7}}
    \proofstepwide[2,3]{\eta_R(m+\Succ(n))}{=}{\eta_R(\Succ(m+n))}{%
      \FormulaRefAuto{PeanoAddSuccRight}[thm]{2,3}}
    \proofstepwide[1,2,3]{}{=}{\eta_R(m+n)\oplus1_R}{%
      \FormulaRefAuto{NaturalSemiringMapSuccessor}[thm]{1,5}}
    \proofstepwide[1,2,3,4]{}{=}{%
      \bigl(\eta_R(m)\oplus\eta_R(n)\bigr)\oplus1_R}{\rIE{4}}
    \proofstepwide[1,2,3]{}{=}{%
      \eta_R(m)\oplus\bigl(\eta_R(n)\oplus1_R\bigr)}{\rIE{6}}
    \proofstepwide[1,2,3]{}{=}{\eta_R(m)\oplus\eta_R(\Succ(n))}{\rIE{8}}
    \proofstepwidestar[1,2,3,4]{%
      \eta_R(m+\Succ(n))=\eta_R(m)\oplus\eta_R(\Succ(n))}{%
      \rChain{9,13}}
  \closeproofpartwideindR
\end{tabproofsplitwide}

\FormulaThmDeltaK[Die kanonische Abbildung erhält die Addition]{%
  \begin{aligned}[t]
    &\SemiringAlg{R}{\oplus}{\odot}{0_R}{1_R}\dsep m,n\in\mathbb N\\
    &\vdash
      \NaturalSemiringMap{R,\oplus,\odot,0_R,1_R}(m+n)
      =
      \NaturalSemiringMap{R,\oplus,\odot,0_R,1_R}(m)
      \oplus
      \NaturalSemiringMap{R,\oplus,\odot,0_R,1_R}(n)
  \end{aligned}
}{NaturalSemiringMapAddition}{%
  \DeltaRow{Mengen}{R\dsep0_R\dsep1_R\dsep m\dsep n\dsep k}
  \DeltaRow{Funktionen}{\oplus\dsep\odot}
}
\begin{tabproofwide}
  \proofstepwidestar[1]{%
    \SemiringAlg{R}{\oplus}{\odot}{0_R}{1_R}}{\rA}
  \proofstepwidestar[2]{m\in\mathbb N}{\rA}
  \proofstepwidestar[3]{n\in\mathbb N}{\rA}
  \proofstepwidestar[1,2]{\eta_R(m+0)=\eta_R(m)\oplus\eta_R(0)}{%
    \FormulaRefAuto{NaturalSemiringMapAdditionZero}[thm]{1,2}}
  \proofstepwidestar[5]{k\in\mathbb N}{\rA}
  \proofstepwidestar[6]{\eta_R(m+k)=\eta_R(m)\oplus\eta_R(k)}{\rA}
  \proofstepwidestar[1,2,5,6]{%
    \eta_R(m+\Succ(k))=\eta_R(m)\oplus\eta_R(\Succ(k))}{%
    \FormulaRefAuto{NaturalSemiringMapAdditionStep}[thm]{1,2,5,6}}
  \proofstepwidestar[1,2,3]{\eta_R(m+n)=\eta_R(m)\oplus\eta_R(n)}{%
    \FormulaRefAuto{PeanoInductionRule}[ax]{4,7,3}}
\end{tabproofwide}

\FormulaThmDeltaK[Nullfall der Multiplikationserhaltung]{%
  \begin{aligned}[t]
    &\SemiringAlg{R}{\oplus}{\odot}{0_R}{1_R}\dsep m\in\mathbb N\\
    &\vdash
      \NaturalSemiringMap{R,\oplus,\odot,0_R,1_R}(m\cdot0)
      =
      \NaturalSemiringMap{R,\oplus,\odot,0_R,1_R}(m)
      \odot
      \NaturalSemiringMap{R,\oplus,\odot,0_R,1_R}(0)
  \end{aligned}
}{NaturalSemiringMapMultiplicationZero}{%
  \DeltaRow{Mengen}{R\dsep0_R\dsep1_R\dsep m}
  \DeltaRow{Funktionen}{\oplus\dsep\odot}
}
\begin{tabproofwide}
  \proofstepwidestar[1]{%
    \SemiringAlg{R}{\oplus}{\odot}{0_R}{1_R}}{\rA}
  \proofstepwidestar[2]{m\in\mathbb N}{\rA}
  \proofstepwidestar[1]{\eta_R\colon\mathbb N\to R}{%
    \FormulaRefAuto{NaturalSemiringMapFunction}[thm]{1}}
  \proofstepwidestar[1,2]{\eta_R(m)\in R}{%
    \FormulaRefAuto{F\colon A\to B\dsep a\in A\vdash F(a)\in B}{3,2}}
  \proofstepwidestar[1,2]{\eta_R(m)\odot0_R=0_R}{%
    \FormulaRefAuto{SemiringRightZeroAxiom}[ax]{1,4}}
  \proofstepwidestar[1,2]{0_R=\eta_R(m)\odot0_R}{%
    \FormulaRefAuto{a=b\vdash b=a}{5}}
  \proofstepwidestar[1]{\eta_R(0)=0_R}{%
    \FormulaRefAuto{NaturalSemiringMapZero}[thm]{1}}
  \proofstepwidestar[1]{0_R=\eta_R(0)}{%
    \FormulaRefAuto{a=b\vdash b=a}{7}}
  \proofstepwide[2]{\eta_R(m\cdot0)}{=}{\eta_R(0)}{%
    \FormulaRefAuto{PeanoMulRightZero}[thm]{2}}
  \proofstepwide[1,2]{}{=}{0_R}{\rIE{7}}
  \proofstepwide[1,2]{}{=}{\eta_R(m)\odot0_R}{\rIE{6}}
  \proofstepwide[1,2]{}{=}{\eta_R(m)\odot\eta_R(0)}{\rIE{8}}
  \proofstepwidestar[1,2]{%
    \eta_R(m\cdot0)=\eta_R(m)\odot\eta_R(0)}{%
    \rChain{9,12}}
\end{tabproofwide}

\FormulaThmDeltaK[Nachfolgerschritt der Multiplikationserhaltung]{%
  \begin{aligned}[t]
    &\SemiringAlg{R}{\oplus}{\odot}{0_R}{1_R}\dsep m,n\in\mathbb N\dsep{}\\[-2pt]
    &\NaturalSemiringMap{R,\oplus,\odot,0_R,1_R}(m\cdot n)
      =
      \NaturalSemiringMap{R,\oplus,\odot,0_R,1_R}(m)
      \odot
      \NaturalSemiringMap{R,\oplus,\odot,0_R,1_R}(n)\\
    &\vdash
      \NaturalSemiringMap{R,\oplus,\odot,0_R,1_R}(m\cdot\Succ(n))
      =
      \NaturalSemiringMap{R,\oplus,\odot,0_R,1_R}(m)
      \odot
      \NaturalSemiringMap{R,\oplus,\odot,0_R,1_R}(\Succ(n))
  \end{aligned}
}{NaturalSemiringMapMultiplicationStep}{%
  \DeltaRow{Mengen}{R\dsep0_R\dsep1_R\dsep m\dsep n}
  \DeltaRow{Funktionen}{\oplus\dsep\odot}
}
\begin{notation*}
\[
  x:=\eta_R(m),
  \qquad
  y:=\eta_R(n).
\]
\end{notation*}
\begin{tabproofsplitwide}
  \proofpartwideindR[Rechtsneutrale Darstellung von \(x\)]{%
    \SemiringAlg{R}{\oplus}{\odot}{0_R}{1_R}\dsep m\in\mathbb N
    \vdash x=x\odot1_R
  }{%
    \SemiringAlg{R}{\oplus}{\odot}{0_R}{1_R}\dsep m\in\mathbb N
    \vdash x=x\odot1_R
  }[NaturalSemiringMapMultiplicationStepRightIdentity]
    \proofstepwidestar[1]{%
      \SemiringAlg{R}{\oplus}{\odot}{0_R}{1_R}}{\rA}
    \proofstepwidestar[2]{m\in\mathbb N}{\rA}
    \proofstepwidestar[1]{\eta_R\colon\mathbb N\to R}{%
      \FormulaRefAuto{NaturalSemiringMapFunction}[thm]{1}}
    \proofstepwidestar[1,2]{x\in R}{%
      \FormulaRefAuto{F\colon A\to B\dsep a\in A\vdash F(a)\in B}{3,2}}
    \proofstepwidestar[1]{\MonoidAlg{R}{\odot}{1_R}}{%
      \FormulaRefAuto{SemiringMultiplicativeMonoidAxiom}[ax]{1}}
    \proofstepwidestar[1,2]{x\odot1_R=x}{%
      \FormulaRefAuto{MonoidRightIdentityAxiom}[ax]{5,4}}
    \proofstepwidestar[1,2]{x=x\odot1_R}{%
      \FormulaRefAuto{a=b\vdash b=a}{6}}
  \closeproofpartwideindR

  \proofpartwideindR[Umgekehrte Distributivitätsgleichung]{%
    \begin{aligned}[t]
      &\SemiringAlg{R}{\oplus}{\odot}{0_R}{1_R}\dsep m,n\in\mathbb N
      \vdash{}\\[-2pt]
      &\bigl(x\odot y\bigr)\oplus\bigl(x\odot1_R\bigr)
       =x\odot\bigl(y\oplus1_R\bigr)
    \end{aligned}
  }{%
    \begin{aligned}[t]
      &\SemiringAlg{R}{\oplus}{\odot}{0_R}{1_R}\dsep m,n\in\mathbb N
      \vdash{}\\[-2pt]
      &\bigl(x\odot y\bigr)\oplus\bigl(x\odot1_R\bigr)
       =x\odot\bigl(y\oplus1_R\bigr)
    \end{aligned}
  }[NaturalSemiringMapMultiplicationStepDistributivity]
    \proofstepwidestar[1]{%
      \SemiringAlg{R}{\oplus}{\odot}{0_R}{1_R}}{\rA}
    \proofstepwidestar[2]{m\in\mathbb N}{\rA}
    \proofstepwidestar[3]{n\in\mathbb N}{\rA}
    \proofstepwidestar[1]{\eta_R\colon\mathbb N\to R}{%
      \FormulaRefAuto{NaturalSemiringMapFunction}[thm]{1}}
    \proofstepwidestar[1,2]{x\in R}{%
      \FormulaRefAuto{F\colon A\to B\dsep a\in A\vdash F(a)\in B}{4,2}}
    \proofstepwidestar[1,3]{y\in R}{%
      \FormulaRefAuto{F\colon A\to B\dsep a\in A\vdash F(a)\in B}{4,3}}
    \proofstepwidestar[1]{\MonoidAlg{R}{\odot}{1_R}}{%
      \FormulaRefAuto{SemiringMultiplicativeMonoidAxiom}[ax]{1}}
    \proofstepwidestar[1]{1_R\in R}{%
      \FormulaRefAuto{MonoidIdentityCarrierAxiom}[ax]{7}}
    \proofstepwidestar[1,2,3]{%
      x\odot\bigl(y\oplus1_R\bigr)
      =\bigl(x\odot y\bigr)\oplus\bigl(x\odot1_R\bigr)}{%
      \FormulaRefAuto{SemiringLeftDistributiveAxiom}[ax]{1,5,6,8}}
    \proofstepwidestar[1,2,3]{%
      \bigl(x\odot y\bigr)\oplus\bigl(x\odot1_R\bigr)
      =x\odot\bigl(y\oplus1_R\bigr)}{%
      \FormulaRefAuto{a=b\vdash b=a}{9}}
  \closeproofpartwideindR

  \proofpartwideindR[Nachfolgerdarstellung von \(y\)]{%
    \SemiringAlg{R}{\oplus}{\odot}{0_R}{1_R}\dsep n\in\mathbb N
    \vdash y\oplus1_R=\eta_R(\Succ(n))
  }{%
    \SemiringAlg{R}{\oplus}{\odot}{0_R}{1_R}\dsep n\in\mathbb N
    \vdash y\oplus1_R=\eta_R(\Succ(n))
  }[NaturalSemiringMapMultiplicationStepSuccessor]
    \proofstepwidestar[1]{%
      \SemiringAlg{R}{\oplus}{\odot}{0_R}{1_R}}{\rA}
    \proofstepwidestar[2]{n\in\mathbb N}{\rA}
    \proofstepwidestar[1,2]{\eta_R(\Succ(n))=y\oplus1_R}{%
      \FormulaRefAuto{NaturalSemiringMapSuccessor}[thm]{1,2}}
    \proofstepwidestar[1,2]{y\oplus1_R=\eta_R(\Succ(n))}{%
      \FormulaRefAuto{a=b\vdash b=a}{3}}
  \closeproofpartwideindR

  \proofpartwide{Gleichheitskette}
    \proofstepwidestar[1]{%
      \SemiringAlg{R}{\oplus}{\odot}{0_R}{1_R}}{\rA}
    \proofstepwidestar[2]{m\in\mathbb N}{\rA}
    \proofstepwidestar[3]{n\in\mathbb N}{\rA}
    \proofstepwidestar[4]{\eta_R(m\cdot n)=x\odot y}{\rA}
    \proofstepwidestar[2,3]{m\cdot n\in\mathbb N}{%
      \FormulaRefAuto{PeanoMulClosure}[thm]{2,3}}
    \proofstepwidestar[1,2]{x=x\odot1_R}{%
      \FormulaRefAuto{NaturalSemiringMapMultiplicationStepRightIdentity}[thm]{1,2}}
    \proofstepwidestar[1,2,3]{%
      \bigl(x\odot y\bigr)\oplus\bigl(x\odot1_R\bigr)
      =x\odot\bigl(y\oplus1_R\bigr)}{%
      \FormulaRefAuto{NaturalSemiringMapMultiplicationStepDistributivity}[thm]{1,2,3}}
    \proofstepwidestar[1,3]{y\oplus1_R=\eta_R(\Succ(n))}{%
      \FormulaRefAuto{NaturalSemiringMapMultiplicationStepSuccessor}[thm]{1,3}}
    \proofstepwide[2,3]{\eta_R(m\cdot\Succ(n))}{=}{\eta_R(m\cdot n+m)}{%
      \FormulaRefAuto{PeanoMulSuccRight}[thm]{2,3}}
    \proofstepwide[1,2,3]{}{=}{\eta_R(m\cdot n)\oplus x}{%
      \FormulaRefAuto{NaturalSemiringMapAddition}[thm]{1,5,2}}
    \proofstepwide[1,2,3,4]{}{=}{%
      \bigl(x\odot y\bigr)\oplus x}{\rIE{4}}
    \proofstepwide[1,2,3,4]{}{=}{%
      \bigl(x\odot y\bigr)\oplus\bigl(x\odot1_R\bigr)}{\rIE{6}}
    \proofstepwide[1,2,3,4]{}{=}{%
      x\odot\bigl(y\oplus1_R\bigr)}{\rIE{7}}
    \proofstepwide[1,2,3,4]{}{=}{x\odot\eta_R(\Succ(n))}{\rIE{8}}
    \proofstepwidestar[1,2,3,4]{%
      \eta_R(m\cdot\Succ(n))=x\odot\eta_R(\Succ(n))}{%
      \rChain{9,14}}
  \closeproofpartwide
\end{tabproofsplitwide}

\FormulaThmDeltaK[Die kanonische Abbildung erhält die Multiplikation]{%
  \begin{aligned}[t]
    &\SemiringAlg{R}{\oplus}{\odot}{0_R}{1_R}\dsep m,n\in\mathbb N\\
    &\vdash
      \NaturalSemiringMap{R,\oplus,\odot,0_R,1_R}(m\cdot n)
      =
      \NaturalSemiringMap{R,\oplus,\odot,0_R,1_R}(m)
      \odot
      \NaturalSemiringMap{R,\oplus,\odot,0_R,1_R}(n)
  \end{aligned}
}{NaturalSemiringMapMultiplication}{%
  \DeltaRow{Mengen}{R\dsep0_R\dsep1_R\dsep k\dsep m\dsep n}
  \DeltaRow{Funktionen}{\oplus\dsep\odot}
}
\begin{tabproofwide}
  \proofstepwidestar[1]{%
    \SemiringAlg{R}{\oplus}{\odot}{0_R}{1_R}}{\rA}
  \proofstepwidestar[2]{m\in\mathbb N}{\rA}
  \proofstepwidestar[3]{n\in\mathbb N}{\rA}
  \proofstepwidestar[1,2]{\eta_R(m\cdot0)=\eta_R(m)\odot\eta_R(0)}{%
    \FormulaRefAuto{NaturalSemiringMapMultiplicationZero}[thm]{1,2}}
  \proofstepwidestar[5]{k\in\mathbb N}{\rA}
  \proofstepwidestar[6]{\eta_R(m\cdot k)=\eta_R(m)\odot\eta_R(k)}{\rA}
  \proofstepwidestar[1,2,5,6]{%
    \eta_R(m\cdot\Succ(k))=\eta_R(m)\odot\eta_R(\Succ(k))}{%
    \FormulaRefAuto{NaturalSemiringMapMultiplicationStep}[thm]{1,2,5,6}}
  \proofstepwidestar[1,2,3]{\eta_R(m\cdot n)=\eta_R(m)\odot\eta_R(n)}{%
    \FormulaRefAuto{PeanoInductionRule}[ax]{4,7,3}}
\end{tabproofwide}

\begin{notation*}
  \[
    \eta_R\coloneqq
    \NaturalSemiringMap{R,\oplus,\odot,0_R,1_R}.
  \]
\end{notation*}

\FormulaThmDeltaK[Additiver Monoidanteil der kanonischen Abbildung]{%
  \begin{aligned}[t]
    &\SemiringAlg{R}{\oplus}{\odot}{0_R}{1_R}\vdash{}\\[-2pt]
    &\MonoidHom{\NaturalSemiringMap{R,\oplus,\odot,0_R,1_R}}
      {\mathbb N,+,0}{R,\oplus,0_R}
  \end{aligned}
}{NaturalSemiringCanonicalAdditiveMonoidHom}{%
  \DeltaRow{Mengen}{R\dsep0_R\dsep1_R\dsep m\dsep n}
  \DeltaRow{Funktionen}{\oplus\dsep\odot}
}
\begin{tabproofwide}
    \proofstepwidestar[1]{%
      \SemiringAlg{R}{\oplus}{\odot}{0_R}{1_R}}{\rA}
    \proofstepwidestar[]{\MonoidAlg{\mathbb N}{+}{0}}{%
      \FormulaRefAuto{NaturalAdditionMonoid}[thm]}
    \proofstepwidestar[1]{\MonoidAlg{R}{\oplus}{0_R}}{%
      \FormulaRefAuto{CommutativeMonoidBaseMonoidAxiom}[ax]{%
        \FormulaRefAuto{SemiringAdditiveCommutativeMonoidAxiom}[ax]{1}}}
    \proofstepwidestar[]{\SemigroupAlg{\mathbb N}{+}}{%
      \FormulaRefAuto{MonoidBaseSemigroupAxiom}[ax]{2}}
    \proofstepwidestar[1]{\SemigroupAlg{R}{\oplus}}{%
      \FormulaRefAuto{MonoidBaseSemigroupAxiom}[ax]{3}}
    \proofstepwidestar[1]{\eta_R\colon\mathbb N\to R}{%
      \FormulaRefAuto{NaturalSemiringMapFunction}[thm]{1}}
    \proofstepwidestar[7]{m\in\mathbb N}{\rA}
    \proofstepwidestar[8]{n\in\mathbb N}{\rA}
    \proofstepwidestar[1,7,8]{%
      \eta_R(m+n)=\eta_R(m)\oplus\eta_R(n)}{%
      \FormulaRefAuto{NaturalSemiringMapAddition}[thm]{1,7,8}}
    \proofstepwidestar[1]{%
      \SemigroupHom{\eta_R}{\mathbb N,+}{R,\oplus}}{%
      \FormulaRefAuto{SemigroupHomDef}[def]{4,5,6,9}}
    \proofstepwidestar[1]{\eta_R(0)=0_R}{%
      \FormulaRefAuto{NaturalSemiringMapZero}[thm]{1}}
    \proofstepwidestar[1]{%
      \MonoidHom{\eta_R}{\mathbb N,+,0}{R,\oplus,0_R}}{%
      \FormulaRefAuto{MonoidHomDef}[def]{2,3,10,11}}
\end{tabproofwide}

\FormulaThmDeltaK[Multiplikativer Monoidanteil der kanonischen Abbildung]{%
  \begin{aligned}[t]
    &\SemiringAlg{R}{\oplus}{\odot}{0_R}{1_R}\vdash{}\\[-2pt]
    &\MonoidHom{\NaturalSemiringMap{R,\oplus,\odot,0_R,1_R}}
      {\mathbb N,\cdot,1}{R,\odot,1_R}
  \end{aligned}
}{NaturalSemiringCanonicalMultiplicativeMonoidHom}{%
  \DeltaRow{Mengen}{R\dsep0_R\dsep1_R\dsep m\dsep n}
  \DeltaRow{Funktionen}{\oplus\dsep\odot}
}
\begin{tabproofwide}
    \proofstepwidestar[1]{%
      \SemiringAlg{R}{\oplus}{\odot}{0_R}{1_R}}{\rA}
    \proofstepwidestar[]{\MonoidAlg{\mathbb N}{\cdot}{1}}{%
      \FormulaRefAuto{NaturalMultiplicationMonoid}[thm]}
    \proofstepwidestar[1]{\MonoidAlg{R}{\odot}{1_R}}{%
      \FormulaRefAuto{SemiringMultiplicativeMonoidAxiom}[ax]{1}}
    \proofstepwidestar[]{\SemigroupAlg{\mathbb N}{\cdot}}{%
      \FormulaRefAuto{MonoidBaseSemigroupAxiom}[ax]{2}}
    \proofstepwidestar[1]{\SemigroupAlg{R}{\odot}}{%
      \FormulaRefAuto{MonoidBaseSemigroupAxiom}[ax]{3}}
    \proofstepwidestar[1]{\eta_R\colon\mathbb N\to R}{%
      \FormulaRefAuto{NaturalSemiringMapFunction}[thm]{1}}
    \proofstepwidestar[7]{m\in\mathbb N}{\rA}
    \proofstepwidestar[8]{n\in\mathbb N}{\rA}
    \proofstepwidestar[1,7,8]{%
      \eta_R(m\cdot n)=\eta_R(m)\odot\eta_R(n)}{%
      \FormulaRefAuto{NaturalSemiringMapMultiplication}[thm]{1,7,8}}
    \proofstepwidestar[1]{%
      \SemigroupHom{\eta_R}{\mathbb N,\cdot}{R,\odot}}{%
      \FormulaRefAuto{SemigroupHomDef}[def]{4,5,6,9}}
    \proofstepwidestar[1]{\eta_R(1)=1_R}{%
      \FormulaRefAuto{NaturalSemiringMapOne}[thm]{1}}
    \proofstepwidestar[1]{%
      \MonoidHom{\eta_R}{\mathbb N,\cdot,1}{R,\odot,1_R}}{%
      \FormulaRefAuto{MonoidHomDef}[def]{2,3,10,11}}
\end{tabproofwide}

\FormulaThmDeltaK[Die kanonische Abbildung ist ein Halbringhomomorphismus]{%
  \SemiringAlg{R}{\oplus}{\odot}{0_R}{1_R}
  \vdash
  \SemiringHom{\NaturalSemiringMap{R,\oplus,\odot,0_R,1_R}}
    {\mathbb N,+,\cdot,0,1}{R,\oplus,\odot,0_R,1_R}%
}{NaturalSemiringCanonicalHomomorphism}{%
  \DeltaRow{Mengen}{R\dsep0_R\dsep1_R}
  \DeltaRow{Funktionen}{\oplus\dsep\odot}
}
\begin{tabproofwide}
    \proofstepwidestar[1]{%
      \SemiringAlg{R}{\oplus}{\odot}{0_R}{1_R}}{\rA}
    \proofstepwidestar[]{%
      \SemiringAlg{\mathbb N}{+}{\cdot}{0}{1}}{%
      \FormulaRefAuto{NaturalUnitalSemiring}[thm]}
    \proofstepwidestar[1]{%
      \MonoidHom{\NaturalSemiringMap{R,\oplus,\odot,0_R,1_R}}
        {\mathbb N,+,0}{R,\oplus,0_R}}{%
      \FormulaRefAuto{NaturalSemiringCanonicalAdditiveMonoidHom}[thm]{1}}
    \proofstepwidestar[1]{%
      \MonoidHom{\NaturalSemiringMap{R,\oplus,\odot,0_R,1_R}}
        {\mathbb N,\cdot,1}{R,\odot,1_R}}{%
      \FormulaRefAuto{NaturalSemiringCanonicalMultiplicativeMonoidHom}[thm]{1}}
    \proofstepwidestar[1]{%
      \SemiringHom{\NaturalSemiringMap{R,\oplus,\odot,0_R,1_R}}
        {\mathbb N,+,\cdot,0,1}{R,\oplus,\odot,0_R,1_R}}{%
      \FormulaRefAuto{SemiringHomDef}[def]{2,1,3,4}}
\end{tabproofwide}

\FormulaThmDeltaK[Eindeutigkeit auf jedem natürlichen Argument]{%
  \begin{aligned}[t]
    &\SemiringAlg{R}{\oplus}{\odot}{0_R}{1_R}\dsep
      \SemiringHom{f}{\mathbb N,+,\cdot,0,1}
        {R,\oplus,\odot,0_R,1_R}\dsep n\in\mathbb N\\
    &\vdash
      f(n)=\NaturalSemiringMap{R,\oplus,\odot,0_R,1_R}(n)
  \end{aligned}
}{NaturalSemiringHomPointwiseUnique}{%
  \DeltaRow{Mengen}{R\dsep0_R\dsep1_R\dsep k\dsep n}
  \DeltaRow{Funktionen}{f\dsep\oplus\dsep\odot}
}
\begin{tabproofsplitwide}
  \proofpartwideindR[Induktionsanfang]{%
    \begin{aligned}[t]
      &\SemiringAlg{R}{\oplus}{\odot}{0_R}{1_R}\dsep
        \SemiringHom{f}{\mathbb N,+,\cdot,0,1}
          {R,\oplus,\odot,0_R,1_R}\vdash{}\\[-2pt]
      &f(0)=\eta_R(0)
    \end{aligned}
  }{%
    \begin{aligned}[t]
      &\SemiringAlg{R}{\oplus}{\odot}{0_R}{1_R}\dsep
        \SemiringHom{f}{\mathbb N,+,\cdot,0,1}
          {R,\oplus,\odot,0_R,1_R}\vdash{}\\[-2pt]
      &f(0)=\eta_R(0)
    \end{aligned}
  }[NaturalSemiringHomPointwiseUniqueZero]
    \proofstepwidestar[1]{%
      \SemiringAlg{R}{\oplus}{\odot}{0_R}{1_R}}{\rA}
    \proofstepwidestar[2]{%
      \SemiringHom{f}{\mathbb N,+,\cdot,0,1}
        {R,\oplus,\odot,0_R,1_R}}{\rA}
    \proofstepwidestar[2]{f(0)=0_R\land f(1)=1_R}{%
      \FormulaRefAuto{SemiringHomConstantsAxiom}[ax]{2}}
    \proofstepwidestar[1]{\eta_R(0)=0_R}{%
      \FormulaRefAuto{NaturalSemiringMapZero}[thm]{1}}
    \proofstepwidestar[1]{0_R=\eta_R(0)}{%
      \FormulaRefAuto{a=b\vdash b=a}{4}}
    \proofstepwide[2]{f(0)}{=}{0_R}{\rAEa{3}}
    \proofstepwide[1,2]{}{=}{\eta_R(0)}{\rIE{5}}
    \proofstepwidestar[1,2]{f(0)=\eta_R(0)}{\rChain{6,7}}
  \closeproofpartwideindR

  \proofpartwideindR[Induktionsschritt]{%
    \begin{aligned}[t]
      &\SemiringAlg{R}{\oplus}{\odot}{0_R}{1_R}\dsep
        \SemiringHom{f}{\mathbb N,+,\cdot,0,1}
          {R,\oplus,\odot,0_R,1_R}\dsep n\in\mathbb N\dsep{}\\[-2pt]
      &f(n)=\eta_R(n)\vdash f(\Succ(n))=\eta_R(\Succ(n))
    \end{aligned}
  }{%
    \begin{aligned}[t]
      &\SemiringAlg{R}{\oplus}{\odot}{0_R}{1_R}\dsep
        \SemiringHom{f}{\mathbb N,+,\cdot,0,1}
          {R,\oplus,\odot,0_R,1_R}\dsep n\in\mathbb N\dsep{}\\[-2pt]
      &f(n)=\eta_R(n)\vdash f(\Succ(n))=\eta_R(\Succ(n))
    \end{aligned}
  }[NaturalSemiringHomPointwiseUniqueStep]
    \proofstepwidestar[1]{%
      \SemiringAlg{R}{\oplus}{\odot}{0_R}{1_R}}{\rA}
    \proofstepwidestar[2]{%
      \SemiringHom{f}{\mathbb N,+,\cdot,0,1}
        {R,\oplus,\odot,0_R,1_R}}{\rA}
    \proofstepwidestar[3]{n\in\mathbb N}{\rA}
    \proofstepwidestar[4]{f(n)=\eta_R(n)}{\rA}
    \proofstepwidestar[2]{f(0)=0_R\land f(1)=1_R}{%
      \FormulaRefAuto{SemiringHomConstantsAxiom}[ax]{2}}
    \proofstepwidestar[3]{n+1=\Succ(n)}{%
      \FormulaRefAuto{PeanoAddOneSucc}[thm]{3}}
    \proofstepwidestar[3]{\Succ(n)=n+1}{%
      \FormulaRefAuto{a=b\vdash b=a}{6}}
    \proofstepwidestar[]{1\in\mathbb N}{%
      \FormulaRefAuto{PeanoOneInN}[thm]}
    \proofstepwidestar[2]{f(1)=1_R}{\rAEb{5}}
    \proofstepwidestar[1,3]{\eta_R(\Succ(n))=\eta_R(n)\oplus1_R}{%
      \FormulaRefAuto{NaturalSemiringMapSuccessor}[thm]{1,3}}
    \proofstepwidestar[1,3]{\eta_R(n)\oplus1_R=\eta_R(\Succ(n))}{%
      \FormulaRefAuto{a=b\vdash b=a}{10}}
    \proofstepwide[3]{f(\Succ(n))}{=}{f(n+1)}{\rIE{7}}
    \proofstepwide[2,3]{}{=}{f(n)\oplus f(1)}{%
      \FormulaRefAuto{SemiringHomAdditionAxiom}[ax]{2,3,8}}
    \proofstepwide[1,2,3,4]{}{=}{\eta_R(n)\oplus f(1)}{\rIE{4}}
    \proofstepwide[1,2,3,4]{}{=}{\eta_R(n)\oplus1_R}{\rIE{9}}
    \proofstepwide[1,2,3,4]{}{=}{\eta_R(\Succ(n))}{\rIE{11}}
    \proofstepwidestar[1,2,3,4]{%
      f(\Succ(n))=\eta_R(\Succ(n))}{\rChain{12,16}}
  \closeproofpartwideindR

  \proofpartwide{Induktionsschluss}
    \proofstepwidestar[1]{%
      \SemiringAlg{R}{\oplus}{\odot}{0_R}{1_R}}{\rA}
    \proofstepwidestar[2]{%
      \SemiringHom{f}{\mathbb N,+,\cdot,0,1}
        {R,\oplus,\odot,0_R,1_R}}{\rA}
    \proofstepwidestar[3]{n\in\mathbb N}{\rA}
    \proofstepwidestar[1,2]{f(0)=\eta_R(0)}{%
      \FormulaRefAuto{NaturalSemiringHomPointwiseUniqueZero}[thm]{1,2}}
    \proofstepwidestar[5]{k\in\mathbb N}{\rA}
    \proofstepwidestar[6]{f(k)=\eta_R(k)}{\rA}
    \proofstepwidestar[1,2,5,6]{f(\Succ(k))=\eta_R(\Succ(k))}{%
      \FormulaRefAuto{NaturalSemiringHomPointwiseUniqueStep}[thm]{1,2,5,6}}
    \proofstepwidestar[1,2,3]{f(n)=\eta_R(n)}{%
      \FormulaRefAuto{PeanoInductionRule}[ax]{4,7,3}}
  \closeproofpartwide
\end{tabproofsplitwide}

\FormulaThmDeltaK[Universelle Eigenschaft der natürlichen Zahlen]{%
  \begin{aligned}[t]
    &\SemiringAlg{R}{\oplus}{\odot}{0_R}{1_R}\vdash{}\\[-2pt]
    &\exists! f\,
      \SemiringHom{f}{\mathbb N,+,\cdot,0,1}
        {R,\oplus,\odot,0_R,1_R}
  \end{aligned}
}{NaturalSemiringUniversalHomomorphism}{%
  \DeltaRow{Mengen}{R\dsep0_R\dsep1_R\dsep n}
  \DeltaRow{Funktionen}{f\dsep\oplus\dsep\odot}
}
\begin{tabproofwide}
  \proofstepwidestar[1]{%
    \SemiringAlg{R}{\oplus}{\odot}{0_R}{1_R}}{\rA}
  \proofstepwidestar[1]{%
    \SemiringHom{\NaturalSemiringMap{R,\oplus,\odot,0_R,1_R}}
      {\mathbb N,+,\cdot,0,1}{R,\oplus,\odot,0_R,1_R}}{%
    \FormulaRefAuto{NaturalSemiringCanonicalHomomorphism}[thm]{1}}
  \proofstepwidestar[3]{%
    \SemiringHom{f}{\mathbb N,+,\cdot,0,1}
      {R,\oplus,\odot,0_R,1_R}}{\rA}
  \proofstepwidestar[3]{f\colon\mathbb N\to R}{%
    \FormulaRefAuto{SemiringHomFunctionAxiom}[ax]{3}}
  \proofstepwidestar[1]{%
    \NaturalSemiringMap{R,\oplus,\odot,0_R,1_R}
    \colon\mathbb N\to R}{%
    \FormulaRefAuto{NaturalSemiringMapFunction}[thm]{1}}
  \proofstepwidestar[6]{n\in\mathbb N}{\rA}
  \proofstepwidestar[1,3,6]{%
    f(n)=\NaturalSemiringMap{R,\oplus,\odot,0_R,1_R}(n)}{%
    \FormulaRefAuto{NaturalSemiringHomPointwiseUnique}[thm]{1,3,6}}
  \proofstepwidestar[1,3]{%
    \forall n\in\mathbb N\,
    \bigl(
      f(n)=\NaturalSemiringMap{R,\oplus,\odot,0_R,1_R}(n)
    \bigr)}{%
    \rUI{\rRI{6,7}}}
  \proofstepwidestar[1,3]{%
    f=\NaturalSemiringMap{R,\oplus,\odot,0_R,1_R}}{%
    \FormulaRefAuto{\forall x\in A\,(F(x)=G(x)) \eqvdash F=G}[thm]{8}}
  \proofstepwidestar[1]{%
    \exists! f\,
    \SemiringHom{f}{\mathbb N,+,\cdot,0,1}
      {R,\oplus,\odot,0_R,1_R}}{%
    \UEI{2,3,9}}
\end{tabproofwide}

\section{Surjektivität und Induktion}

\FormulaThmDeltaK[Surjektivität der kanonischen Abbildung genau bei vollem Bild]{%
  \begin{aligned}[t]
    &\SemiringAlg{R}{\oplus}{\odot}{0_R}{1_R}\\
    &\vdash
      \Bigl(
        \NaturalSemiringMap{R,\oplus,\odot,0_R,1_R}
          \colon\mathbb N\sur R
        \leftrightarrow
        \NaturalSemiringMap{R,\oplus,\odot,0_R,1_R}[\mathbb N]=R
      \Bigr)
  \end{aligned}
}{NaturalSemiringMapSurjectiveIffFullImage}{%
  \DeltaRow{Mengen}{R\dsep0_R\dsep1_R}
  \DeltaRow{Funktionen}{\oplus\dsep\odot}
}
\begin{tabproofwide}
  \proofstepwidestar[1]{%
    \SemiringAlg{R}{\oplus}{\odot}{0_R}{1_R}}{\rA}
  \proofstepwidestar[1]{%
    \NaturalSemiringMap{R,\oplus,\odot,0_R,1_R}
      \colon\mathbb N\to R}{%
    \FormulaRefAuto{NaturalSemiringMapFunction}[thm]{1}}
  \proofstepwidestar[1]{%
    \begin{aligned}[t]
      &\NaturalSemiringMap{R,\oplus,\odot,0_R,1_R}
        \colon\mathbb N\sur R\\
      &\leftrightarrow
        \NaturalSemiringMap{R,\oplus,\odot,0_R,1_R}[\mathbb N]=R
    \end{aligned}
  }{%
    \FormulaRefAuto{SurjectiveIffImageEqualsCodomain}[thm]{2}}
\end{tabproofwide}

\FormulaThmDeltaK[Surjektivität der kanonischen Abbildung genau beim Nachfolgerinduktionsprinzip]{%
  \begin{aligned}[t]
    &\SemiringAlg{R}{\oplus}{\odot}{0_R}{1_R}\\
    &\vdash
      \Bigl(
        \NaturalSemiringMap{R,\oplus,\odot,0_R,1_R}
          \colon\mathbb N\sur R
        \leftrightarrow{}\\[-2pt]
    &\qquad
        \forall A\in\powerset(R)\,
        \Bigl(
          \bigl(
            0_R\in A\land
            \forall x\in A\,
            \bigl(
              \SemiringSuccessor{R,\oplus,\odot,0_R,1_R}(x)\in A
            \bigr)
          \bigr)
          \rightarrow A=R
        \Bigr)
      \Bigr)
  \end{aligned}
}{NaturalSemiringMapSurjectiveIffSuccessorInduction}{%
  \DeltaRow{Mengen}{R\dsep A\dsep0_R\dsep1_R\dsep x}
  \DeltaRow{Funktionen}{\oplus\dsep\odot}
}
\begin{notation*}
\[
  \eta_R:=\NaturalSemiringMap{R,\oplus,\odot,0_R,1_R},
  \qquad
  s_R:=\SemiringSuccessor{R,\oplus,\odot,0_R,1_R}.
\]
\end{notation*}
\begin{tabproofwide}
  \proofstepwidestar[1]{%
    \SemiringAlg{R}{\oplus}{\odot}{0_R}{1_R}}{\rA}
  \proofstepwidestar[1]{0_R\in R}{%
    \FormulaRefAuto{SemiringZeroCarrier}[thm]{1}}
  \proofstepwidestar[1]{%
    s_R\colon R\to R}{%
    \FormulaRefAuto{SemiringSuccessorFunction}[thm]{1}}
  \proofstepwidestar[1]{%
    \begin{aligned}[t]
      &\RecFun{0_R}{s_R}\colon\mathbb N\sur R
      \leftrightarrow{}\\[-2pt]
      &\forall A\in\powerset(R)\,
      \Bigl(
        \bigl(
          0_R\in A\land
          \forall x\in A\,
          \bigl(
            s_R(x)\in A
          \bigr)
        \bigr)
        \rightarrow A=R
      \Bigr)
    \end{aligned}
  }{%
    \FormulaRefAuto{RecFunSurjectiveIffInductionPrinciple}[thm]{2,3}}
  \proofstepwidestar[1]{%
    \eta_R=\RecFun{0_R}{s_R}}{%
    \FormulaRefAuto{NaturalSemiringMapDef}[def]{1}}
  \proofstepwidestar[1]{%
    \begin{aligned}[t]
      &\eta_R\colon\mathbb N\sur R
      \leftrightarrow{}\\[-2pt]
      &\forall A\in\powerset(R)\,
      \Bigl(
        \bigl(
          0_R\in A\land
          \forall x\in A\,
          \bigl(
            s_R(x)\in A
          \bigr)
        \bigr)
        \rightarrow A=R
      \Bigr)
    \end{aligned}
  }{\rIE{5,4}}
\end{tabproofwide}

\section{Injektivität}

\FormulaThmDeltaK[Injektivität schließt Kollisionen aus]{%
  \begin{aligned}[t]
    &\SemiringAlg{R}{\oplus}{\odot}{0_R}{1_R}\dsep
      \NaturalSemiringMap{R,\oplus,\odot,0_R,1_R}
        \colon\mathbb N\inj R\\
    &\vdash
      \forall m,n\in\mathbb N\,
      \Bigl(
        \NaturalSemiringMap{R,\oplus,\odot,0_R,1_R}(m)
        =
        \NaturalSemiringMap{R,\oplus,\odot,0_R,1_R}(n)
        \rightarrow m=n
      \Bigr)
  \end{aligned}
}{NaturalSemiringMapInjectiveNoCollision}{%
  \DeltaRow{Mengen}{R\dsep0_R\dsep1_R\dsep m\dsep n}
  \DeltaRow{Funktionen}{\oplus\dsep\odot}
}
\begin{notation*}
\[
  \eta_R:=\NaturalSemiringMap{R,\oplus,\odot,0_R,1_R}.
\]
\end{notation*}
\begin{tabproofwide}
  \proofstepwidestar[1]{%
    \SemiringAlg{R}{\oplus}{\odot}{0_R}{1_R}}{\rA}
  \proofstepwidestar[2]{%
    \eta_R\colon\mathbb N\inj R}{\rA}
  \proofstepwidestar[3]{m\in\mathbb N}{\rA}
  \proofstepwidestar[4]{n\in\mathbb N}{\rA}
  \proofstepwidestar[5]{%
    \eta_R(m)=\eta_R(n)}{\rA}
  \proofstepwidestar[2,3,4,5]{m=n}{%
    \FormulaRefAuto{Injektivität}[ax]{2,3,4,5}}
  \proofstepwidestar[2,3,4]{%
    \begin{aligned}[t]
      &\eta_R(m)=\eta_R(n)\\
      &\rightarrow m=n
    \end{aligned}
  }{\rRI{5,6}}
  \proofstepwidestar[2,3]{%
    \forall n\in\mathbb N\,
    \Bigl(
      \eta_R(m)=\eta_R(n)
      \rightarrow m=n
    \Bigr)}{\rUI{\rRI{4,7}}}
  \proofstepwidestar[2]{%
    \forall m,n\in\mathbb N\,
    \Bigl(
      \eta_R(m)=\eta_R(n)
      \rightarrow m=n
    \Bigr)}{\rUI{\rRI{3,8}}}
\end{tabproofwide}

\FormulaThmDeltaK[Kollisionsfreiheit liefert Injektivität]{%
  \begin{aligned}[t]
    &\SemiringAlg{R}{\oplus}{\odot}{0_R}{1_R}\dsep{}\\[-2pt]
    &\forall m,n\in\mathbb N\,
      \Bigl(
        \NaturalSemiringMap{R,\oplus,\odot,0_R,1_R}(m)
        =
        \NaturalSemiringMap{R,\oplus,\odot,0_R,1_R}(n)
        \rightarrow m=n
      \Bigr)\\
    &\vdash
      \NaturalSemiringMap{R,\oplus,\odot,0_R,1_R}
        \colon\mathbb N\inj R
  \end{aligned}
}{NaturalSemiringMapNoCollisionInjective}{%
  \DeltaRow{Mengen}{R\dsep0_R\dsep1_R\dsep m\dsep n}
  \DeltaRow{Funktionen}{\oplus\dsep\odot}
}
\begin{tabproofwide}
  \proofstepwidestar[1]{%
    \SemiringAlg{R}{\oplus}{\odot}{0_R}{1_R}}{\rA}
  \proofstepwidestar[2]{%
    \forall m,n\in\mathbb N\,
    \Bigl(
      \NaturalSemiringMap{R,\oplus,\odot,0_R,1_R}(m)
      =
      \NaturalSemiringMap{R,\oplus,\odot,0_R,1_R}(n)
      \rightarrow m=n
    \Bigr)}{\rA}
  \proofstepwidestar[1]{%
    \NaturalSemiringMap{R,\oplus,\odot,0_R,1_R}
      \colon\mathbb N\to R}{%
    \FormulaRefAuto{NaturalSemiringMapFunction}[thm]{1}}
  \proofstepwidestar[1,2]{%
    \NaturalSemiringMap{R,\oplus,\odot,0_R,1_R}
      \colon\mathbb N\inj R}{%
    \FormulaRefAuto{Injektive Funktion}[def]{3,2}}
\end{tabproofwide}

\FormulaThmDeltaK[Injektivität der kanonischen Abbildung genau bei Kollisionsfreiheit]{%
  \begin{aligned}[t]
    &\SemiringAlg{R}{\oplus}{\odot}{0_R}{1_R}\\
    &\vdash
      \Bigl(
        \NaturalSemiringMap{R,\oplus,\odot,0_R,1_R}
          \colon\mathbb N\inj R
        \leftrightarrow{}\\[-2pt]
    &\qquad
        \forall m,n\in\mathbb N\,
        \Bigl(
          \NaturalSemiringMap{R,\oplus,\odot,0_R,1_R}(m)
          =
          \NaturalSemiringMap{R,\oplus,\odot,0_R,1_R}(n)
          \rightarrow m=n
        \Bigr)
      \Bigr)
  \end{aligned}
}{NaturalSemiringMapInjectiveIffNoCollision}{%
  \DeltaRow{Mengen}{R\dsep0_R\dsep1_R\dsep m\dsep n}
  \DeltaRow{Funktionen}{\oplus\dsep\odot}
}
\begin{tabproofwide}
    \proofstepwidestar[1]{%
      \SemiringAlg{R}{\oplus}{\odot}{0_R}{1_R}}{\rA}
    \proofstepwidestar[2]{%
      \NaturalSemiringMap{R,\oplus,\odot,0_R,1_R}
        \colon\mathbb N\inj R}{\rA}
    \proofstepwidestar[1,2]{%
      \forall m,n\in\mathbb N\,
      \Bigl(
        \NaturalSemiringMap{R,\oplus,\odot,0_R,1_R}(m)
        =
        \NaturalSemiringMap{R,\oplus,\odot,0_R,1_R}(n)
        \rightarrow m=n
      \Bigr)}{%
      \FormulaRefAuto{NaturalSemiringMapInjectiveNoCollision}[thm]{1,2}}
    \proofstepwidestar[1]{%
      \begin{aligned}[t]
        &\NaturalSemiringMap{R,\oplus,\odot,0_R,1_R}
          \colon\mathbb N\inj R
        \rightarrow{}\\[-2pt]
        &\forall m,n\in\mathbb N\,
        \Bigl(
          \NaturalSemiringMap{R,\oplus,\odot,0_R,1_R}(m)
          =
          \NaturalSemiringMap{R,\oplus,\odot,0_R,1_R}(n)
          \rightarrow m=n
        \Bigr)
      \end{aligned}
    }{\rRI{2,3}}
    \proofstepwidestar[5]{%
      \forall m,n\in\mathbb N\,
      \Bigl(
        \NaturalSemiringMap{R,\oplus,\odot,0_R,1_R}(m)
        =
        \NaturalSemiringMap{R,\oplus,\odot,0_R,1_R}(n)
        \rightarrow m=n
      \Bigr)}{\rA}
    \proofstepwidestar[1,5]{%
      \NaturalSemiringMap{R,\oplus,\odot,0_R,1_R}
        \colon\mathbb N\inj R}{%
      \FormulaRefAuto{NaturalSemiringMapNoCollisionInjective}[thm]{1,5}}
    \proofstepwidestar[1]{%
      \begin{aligned}[t]
        &\forall m,n\in\mathbb N\,
        \Bigl(
          \NaturalSemiringMap{R,\oplus,\odot,0_R,1_R}(m)
          =
          \NaturalSemiringMap{R,\oplus,\odot,0_R,1_R}(n)
          \rightarrow m=n
        \Bigr)\\
        &\rightarrow
        \NaturalSemiringMap{R,\oplus,\odot,0_R,1_R}
          \colon\mathbb N\inj R
      \end{aligned}
    }{\rRI{5,6}}
    \proofstepwidestar[1]{%
      \begin{aligned}[t]
        &\NaturalSemiringMap{R,\oplus,\odot,0_R,1_R}
          \colon\mathbb N\inj R
        \leftrightarrow{}\\[-2pt]
        &\forall m,n\in\mathbb N\,
        \Bigl(
          \NaturalSemiringMap{R,\oplus,\odot,0_R,1_R}(m)
          =
          \NaturalSemiringMap{R,\oplus,\odot,0_R,1_R}(n)
          \rightarrow m=n
        \Bigr)
      \end{aligned}
    }{\rLRI{4,7}}
\end{tabproofwide}

\FormulaThmDeltaK[Peano-Kriterium für die Injektivität der kanonischen Abbildung]{%
  \begin{aligned}[t]
    &\SemiringAlg{R}{\oplus}{\odot}{0_R}{1_R}\dsep
      \SemiringSuccessor{R,\oplus,\odot,0_R,1_R}\colon R\inj R\dsep{}\\[-2pt]
    &\forall x\in R\,
      \Bigl(
        \SemiringSuccessor{R,\oplus,\odot,0_R,1_R}(x)\neq0_R
      \Bigr)\\
    &\vdash
      \NaturalSemiringMap{R,\oplus,\odot,0_R,1_R}
        \colon\mathbb N\inj R
  \end{aligned}
}{NaturalSemiringMapInjectivePeanoCriterion}{%
  \DeltaRow{Mengen}{R\dsep0_R\dsep1_R\dsep x}
  \DeltaRow{Funktionen}{\oplus\dsep\odot}
}
\begin{tabproofwide}
  \proofstepwidestar[1]{%
    \SemiringAlg{R}{\oplus}{\odot}{0_R}{1_R}}{\rA}
  \proofstepwidestar[2]{%
    \SemiringSuccessor{R,\oplus,\odot,0_R,1_R}\colon R\inj R}{\rA}
  \proofstepwidestar[3]{%
    \forall x\in R\,
    \Bigl(
      \SemiringSuccessor{R,\oplus,\odot,0_R,1_R}(x)\neq0_R
    \Bigr)}{\rA}
  \proofstepwidestar[1]{0_R\in R}{%
    \FormulaRefAuto{SemiringZeroCarrier}[thm]{1}}
  \proofstepwidestar[1,2,3]{%
    \RecFun{0_R}
      {\SemiringSuccessor{R,\oplus,\odot,0_R,1_R}}
      \colon\mathbb N\inj R}{%
      \FormulaRefAuto{%
        m_0\in M\dsep f\colon M\inj M\dsep
        \forall x\in M\,\bigl(f(x)\neq m_0\bigr)
        \vdash \RecFun{m_0}{f}\colon\mathbb{N}\inj M}[thm]{4,2,3}}
  \proofstepwidestar[1]{%
    \NaturalSemiringMap{R,\oplus,\odot,0_R,1_R}
    =
    \RecFun{0_R}
      {\SemiringSuccessor{R,\oplus,\odot,0_R,1_R}}}{%
    \FormulaRefAuto{NaturalSemiringMapDef}[def]{1}}
  \proofstepwidestar[1,2,3]{%
    \NaturalSemiringMap{R,\oplus,\odot,0_R,1_R}
      \colon\mathbb N\inj R}{\rIE{6,5}}
\end{tabproofwide}

\FormulaThmDeltaK[Die kanonische Abbildung in einen endlichen Halbring ist nicht injektiv]{%
  \SemiringAlg{R}{\oplus}{\odot}{0_R}{1_R}\dsep\Finite{R}
  \vdash
  \neg\Bigl(
    \NaturalSemiringMap{R,\oplus,\odot,0_R,1_R}
      \colon\mathbb N\inj R
  \Bigr)
}{NaturalSemiringMapNotInjectiveIntoFiniteSemiring}{%
  \DeltaRow{Mengen}{R\dsep0_R\dsep1_R}
  \DeltaRow{Funktionen}{H\dsep\oplus\dsep\odot}
}
\begin{tabproof}
  \proofstep{1}{%
    \SemiringAlg{R}{\oplus}{\odot}{0_R}{1_R}}{\rA}
  \proofstep{2}{\Finite{R}}{\rA}
  \proofstep{2}{%
    \neg\exists H\colon\mathbb N\inj R}{%
    \FormulaRefAuto{FiniteNoNatInjection}[thm]{2}}
  \proofstep{4}{%
    \NaturalSemiringMap{R,\oplus,\odot,0_R,1_R}
      \colon\mathbb N\inj R}{\rA}
  \proofstep{4}{%
    \exists H\colon\mathbb N\inj R}{\rEI{4}}
  \proofstep{2,4}{\bot}{\rBI{3,5}}
  \proofstep{2}{%
    \neg\Bigl(
      \NaturalSemiringMap{R,\oplus,\odot,0_R,1_R}
        \colon\mathbb N\inj R
    \Bigr)}{\rCI{4,6}}
\end{tabproof}

\section{Isomorphie}

\FormulaThmDeltaK[Isomorphie impliziert Injektivität und Surjektivität der kanonischen Abbildung]{%
  \begin{aligned}[t]
    &\SemiringAlg{R}{\oplus}{\odot}{0_R}{1_R}\dsep
      \SemiringIso{\NaturalSemiringMap{R,\oplus,\odot,0_R,1_R}}
        {\mathbb N,+,\cdot,0,1}
        {R,\oplus,\odot,0_R,1_R}\\
    &\vdash
      \NaturalSemiringMap{R,\oplus,\odot,0_R,1_R}
        \colon\mathbb N\inj R
      \land
      \NaturalSemiringMap{R,\oplus,\odot,0_R,1_R}
        \colon\mathbb N\sur R
  \end{aligned}
}{NaturalSemiringMapIsoInjectiveSurjective}{%
  \DeltaRow{Mengen}{R\dsep0_R\dsep1_R}
  \DeltaRow{Funktionen}{\oplus\dsep\odot}
}
\begin{tabproofwide}
    \proofstepwidestar[1]{%
      \SemiringAlg{R}{\oplus}{\odot}{0_R}{1_R}}{\rA}
    \proofstepwidestar[2]{%
      \SemiringIso{\NaturalSemiringMap{R,\oplus,\odot,0_R,1_R}}
        {\mathbb N,+,\cdot,0,1}
        {R,\oplus,\odot,0_R,1_R}}{\rA}
    \proofstepwidestar[2]{%
      \NaturalSemiringMap{R,\oplus,\odot,0_R,1_R}
        \colon\mathbb N\bij R}{%
      \FormulaRefAuto{SemiringIsoDef}[def]{2}}
    \proofstepwidestar[2]{%
      \NaturalSemiringMap{R,\oplus,\odot,0_R,1_R}
        \colon\mathbb N\inj R}{%
      \FormulaRefAuto{%
        F\colon A\bij B\vdash F\colon A\inj B}[thm]{3}}
    \proofstepwidestar[2]{%
      \NaturalSemiringMap{R,\oplus,\odot,0_R,1_R}
        \colon\mathbb N\sur R}{%
      \FormulaRefAuto{%
        F\colon A\bij B\vdash F\colon A\sur B}[thm]{3}}
    \proofstepwidestar[2]{%
      \NaturalSemiringMap{R,\oplus,\odot,0_R,1_R}
        \colon\mathbb N\inj R
      \land
      \NaturalSemiringMap{R,\oplus,\odot,0_R,1_R}
        \colon\mathbb N\sur R}{\rAI{4,5}}
\end{tabproofwide}

\FormulaThmDeltaK[Injektivität und Surjektivität implizieren Isomorphie der kanonischen Abbildung]{%
  \begin{aligned}[t]
    &\SemiringAlg{R}{\oplus}{\odot}{0_R}{1_R}\dsep{}\\[-2pt]
    &\NaturalSemiringMap{R,\oplus,\odot,0_R,1_R}
      \colon\mathbb N\inj R
    \land
    \NaturalSemiringMap{R,\oplus,\odot,0_R,1_R}
      \colon\mathbb N\sur R\\
    &\vdash
      \SemiringIso{\NaturalSemiringMap{R,\oplus,\odot,0_R,1_R}}
        {\mathbb N,+,\cdot,0,1}
        {R,\oplus,\odot,0_R,1_R}
  \end{aligned}
}{NaturalSemiringMapInjectiveSurjectiveIso}{%
  \DeltaRow{Mengen}{R\dsep0_R\dsep1_R}
  \DeltaRow{Funktionen}{\oplus\dsep\odot}
}
\begin{tabproofwide}
    \proofstepwidestar[1]{%
      \SemiringAlg{R}{\oplus}{\odot}{0_R}{1_R}}{\rA}
    \proofstepwidestar[2]{%
      \NaturalSemiringMap{R,\oplus,\odot,0_R,1_R}
        \colon\mathbb N\inj R
      \land
      \NaturalSemiringMap{R,\oplus,\odot,0_R,1_R}
        \colon\mathbb N\sur R}{\rA}
    \proofstepwidestar[2]{%
      \NaturalSemiringMap{R,\oplus,\odot,0_R,1_R}
        \colon\mathbb N\inj R}{\rAEa{2}}
    \proofstepwidestar[2]{%
      \NaturalSemiringMap{R,\oplus,\odot,0_R,1_R}
        \colon\mathbb N\sur R}{\rAEb{2}}
    \proofstepwidestar[2]{%
      \NaturalSemiringMap{R,\oplus,\odot,0_R,1_R}
        \colon\mathbb N\bij R}{%
      \FormulaRefAuto{%
        F\colon A\inj B\dsep F\colon A\sur B
        \vdash F\colon A\bij B}[thm]{3,4}}
    \proofstepwidestar[1]{%
      \SemiringHom{\NaturalSemiringMap{R,\oplus,\odot,0_R,1_R}}
        {\mathbb N,+,\cdot,0,1}
        {R,\oplus,\odot,0_R,1_R}}{%
      \FormulaRefAuto{NaturalSemiringCanonicalHomomorphism}[thm]{1}}
    \proofstepwidestar[1,2]{%
      \SemiringIso{\NaturalSemiringMap{R,\oplus,\odot,0_R,1_R}}
        {\mathbb N,+,\cdot,0,1}
        {R,\oplus,\odot,0_R,1_R}}{%
      \FormulaRefAuto{SemiringIsoDef}[def]{6,5}}
\end{tabproofwide}

\FormulaThmDeltaK[Isomorphie genau bei Injektivität und Surjektivität der kanonischen Abbildung]{%
  \begin{aligned}[t]
    &\SemiringAlg{R}{\oplus}{\odot}{0_R}{1_R}\\
    &\vdash
      \Bigl(
        \SemiringIso{\NaturalSemiringMap{R,\oplus,\odot,0_R,1_R}}
          {\mathbb N,+,\cdot,0,1}
          {R,\oplus,\odot,0_R,1_R}
        \leftrightarrow{}\\[-2pt]
    &\qquad
        \NaturalSemiringMap{R,\oplus,\odot,0_R,1_R}
          \colon\mathbb N\inj R
        \land
        \NaturalSemiringMap{R,\oplus,\odot,0_R,1_R}
          \colon\mathbb N\sur R
      \Bigr)
  \end{aligned}
}{NaturalSemiringMapIsoIffInjectiveAndSurjective}{%
  \DeltaRow{Mengen}{R\dsep0_R\dsep1_R}
  \DeltaRow{Funktionen}{\oplus\dsep\odot}
}
\begin{tabproofwide}
    \proofstepwidestar[1]{%
      \SemiringAlg{R}{\oplus}{\odot}{0_R}{1_R}}{\rA}
    \proofstepwidestar[2]{%
      \SemiringIso{\NaturalSemiringMap{R,\oplus,\odot,0_R,1_R}}
        {\mathbb N,+,\cdot,0,1}
        {R,\oplus,\odot,0_R,1_R}}{\rA}
    \proofstepwidestar[1,2]{%
      \NaturalSemiringMap{R,\oplus,\odot,0_R,1_R}
        \colon\mathbb N\inj R
      \land
      \NaturalSemiringMap{R,\oplus,\odot,0_R,1_R}
        \colon\mathbb N\sur R}{%
      \FormulaRefAuto{NaturalSemiringMapIsoInjectiveSurjective}[thm]{1,2}}
    \proofstepwidestar[1]{%
      \begin{aligned}[t]
        &\SemiringIso{\NaturalSemiringMap{R,\oplus,\odot,0_R,1_R}}
          {\mathbb N,+,\cdot,0,1}
          {R,\oplus,\odot,0_R,1_R}\\
        &\rightarrow
          \NaturalSemiringMap{R,\oplus,\odot,0_R,1_R}
            \colon\mathbb N\inj R
          \land
          \NaturalSemiringMap{R,\oplus,\odot,0_R,1_R}
            \colon\mathbb N\sur R
      \end{aligned}
    }{\rRI{2,3}}
    \proofstepwidestar[5]{%
      \NaturalSemiringMap{R,\oplus,\odot,0_R,1_R}
        \colon\mathbb N\inj R
      \land
      \NaturalSemiringMap{R,\oplus,\odot,0_R,1_R}
        \colon\mathbb N\sur R}{\rA}
    \proofstepwidestar[1,5]{%
      \SemiringIso{\NaturalSemiringMap{R,\oplus,\odot,0_R,1_R}}
        {\mathbb N,+,\cdot,0,1}
        {R,\oplus,\odot,0_R,1_R}}{%
      \FormulaRefAuto{NaturalSemiringMapInjectiveSurjectiveIso}[thm]{1,5}}
    \proofstepwidestar[1]{%
      \begin{aligned}[t]
        &\NaturalSemiringMap{R,\oplus,\odot,0_R,1_R}
          \colon\mathbb N\inj R
        \land
        \NaturalSemiringMap{R,\oplus,\odot,0_R,1_R}
          \colon\mathbb N\sur R\\
        &\rightarrow
          \SemiringIso{\NaturalSemiringMap{R,\oplus,\odot,0_R,1_R}}
            {\mathbb N,+,\cdot,0,1}
            {R,\oplus,\odot,0_R,1_R}
      \end{aligned}
    }{\rRI{5,6}}
    \proofstepwidestar[1]{%
      \begin{aligned}[t]
        &\SemiringIso{\NaturalSemiringMap{R,\oplus,\odot,0_R,1_R}}
          {\mathbb N,+,\cdot,0,1}
          {R,\oplus,\odot,0_R,1_R}
        \leftrightarrow{}\\[-2pt]
        &\NaturalSemiringMap{R,\oplus,\odot,0_R,1_R}
          \colon\mathbb N\inj R
        \land
        \NaturalSemiringMap{R,\oplus,\odot,0_R,1_R}
          \colon\mathbb N\sur R
      \end{aligned}
    }{\rLRI{4,7}}
\end{tabproofwide}

\section{Fehlen additiver Inverser}

\FormulaThmDeltaK[Darstellung als Nachfolger]{%
  n\in\mathbb{N}\vdash1+n=\Succ(n)%
}{NaturalOnePlusAsSuccessor}{%
  \DeltaRow{Mengen}{\mathbb{N}\dsep n}
}
\begin{tabproofwide}
  \proofstepwidestar[1]{n\in\mathbb{N}}{\rA}
  \proofstepwidestar[]{1\in\mathbb{N}}{%
    \FormulaRefAuto{PeanoOneInN}[thm]}
  \proofstepwide[1]{1+n}{=}{n+1}{%
    \FormulaRefAuto{PeanoAddCommutative}[thm]{2,1}}
  \proofstepwide[1]{}{=}{\Succ(n)}{%
    \FormulaRefAuto{PeanoAddOneSucc}[thm]{1}}
  \proofstepwidestar[1]{1+n=\Succ(n)}{%
    \rChain{3,4}}
\end{tabproofwide}

\FormulaThmDeltaK[Widerspruch zum Nullaxiom]{%
  \exists n\in\mathbb{N}\,(1+n=0)\vdash\bot%
}{NaturalAdditionNotGroupContradiction}{%
  \DeltaRow{Mengen}{\mathbb{N}\dsep n}
}
\begin{tabproofwide}
  \proofstepwidestar[1]{%
    \exists n\in\mathbb{N}\,(1+n=0)}{\rA}
  \proofstepwidestar[2]{n\in\mathbb{N}}{\rA}
  \proofstepwidestar[3]{1+n=0}{\rA}
  \proofstepwidestar[2]{1+n=\Succ(n)}{%
    \FormulaRefAuto{NaturalOnePlusAsSuccessor}{2}}
  \proofstepwide[2]{\Succ(n)}{=}{1+n}{%
    \FormulaRefAuto{a=b\vdash b=a}{4}}
  \proofstepwide[2,3]{}{=}{0}{\rIE{3}}
  \proofstepwidestar[2,3]{\Succ(n)=0}{\rChain{5,6}}
  \proofstepwidestar[2]{\Succ(n)\neq0}{%
    \FormulaRefAuto{PeanoSuccNonzero}[ax]{2}}
  \proofstepwidestar[2,3]{\bot}{\rCE{8,7}}
  \proofstepwidestar[1]{\bot}{\rEE{1,2,3,9}}
\end{tabproofwide}

\FormulaThmDeltaK[Die Eins besitzt kein additives Inverses in \(\mathbb{N}\)]{%
  \neg\exists n\in\mathbb{N}\,(1+n=0)%
}{NaturalAdditionNotGroup}{%
  \DeltaRow{Mengen}{\mathbb{N}\dsep n}
}
\begin{tabproof}
  \proofstep{1}{\exists n\in\mathbb{N}\,(1+n=0)}{\rA}
  \proofstep{1}{\bot}{%
    \FormulaRefAuto{NaturalAdditionNotGroupContradiction}[thm]{1}}
  \proofstep{}{\neg\exists n\in\mathbb{N}\,(1+n=0)}{\rCI{1,2}}
\end{tabproof}

\FormulaThmDeltaK[Spezialisierung auf die Eins]{%
  \begin{aligned}[t]
    &\forall a\in\mathbb{N}\,
      \exists b\in\mathbb{N}\,(a+b=0\land b+a=0)\\
    &\vdash\exists b\in\mathbb{N}\,(1+b=0)
  \end{aligned}
}{NaturalRingInverseRequirementAtOne}{%
  \DeltaRow{Mengen}{\mathbb{N}\dsep a\dsep b}
}
\begin{tabproofwide}
  \proofstepwidestar[1]{%
    \forall a\in\mathbb{N}\,
    \exists b\in\mathbb{N}\,(a+b=0\land b+a=0)}{\rA}
  \proofstepwidestar[]{1\in\mathbb{N}}{%
    \FormulaRefAuto{PeanoOneInN}[thm]}
  \proofstepwidestar[1]{%
    \exists b\in\mathbb{N}\,(1+b=0\land b+1=0)}{%
    \rRE{2,\rUE{1}}}
  \proofstepwidestar[4]{b\in\mathbb{N}}{\rA}
  \proofstepwidestar[5]{1+b=0\land b+1=0}{\rA}
  \proofstepwidestar[5]{1+b=0}{\rAEa{5}}
  \proofstepwidestar[4,5]{%
    \exists b\in\mathbb{N}\,(1+b=0)}{\rEI{4,6}}
  \proofstepwidestar[1]{%
    \exists b\in\mathbb{N}\,(1+b=0)}{\rEE{3,4,5,7}}
\end{tabproofwide}

\FormulaThmDeltaK[Fehlschlagen der Inversenforderung]{%
  \neg\forall a\in\mathbb{N}\,
  \exists b\in\mathbb{N}\,(a+b=0\land b+a=0)%
}{NaturalNumbersNotRingConclusion}{%
  \DeltaRow{Mengen}{\mathbb{N}\dsep a\dsep b}
}
\begin{tabproofwide}
  \proofstepwidestar[]{%
    \neg\exists b\in\mathbb{N}\,(1+b=0)}{%
    \FormulaRefAuto{NaturalAdditionNotGroup}[thm]}
  \proofstepwidestar[2]{%
    \forall a\in\mathbb{N}\,
    \exists b\in\mathbb{N}\,(a+b=0\land b+a=0)}{\rA}
  \proofstepwidestar[2]{%
    \exists b\in\mathbb{N}\,(1+b=0)}{%
    \FormulaRefAuto{NaturalRingInverseRequirementAtOne}[thm]{2}}
  \proofstepwidestar[2]{\bot}{\rCE{1,3}}
  \proofstepwidestar[]{%
    \neg\forall a\in\mathbb{N}\,
    \exists b\in\mathbb{N}\,(a+b=0\land b+a=0)}{%
    \rCI{2,4}}
\end{tabproofwide}

\FormulaThmDeltaK[Die natürlichen Zahlen erfüllen nicht die additive Inversenforderung eines Rings]{%
  \begin{aligned}[t]
    &\CommutativeSemiringAlg{\mathbb{N}}{+}{\cdot}{0}{1}\land{}\\[-2pt]
    &\neg\forall a\in\mathbb{N}\,
      \exists b\in\mathbb{N}\,(a+b=0\land b+a=0)
  \end{aligned}
}{NaturalNumbersNotRing}{%
  \DeltaRow{Mengen}{\mathbb{N}\dsep a\dsep b}
}
\begin{tabproofwide}
  \proofstepwidestar[]{%
    \CommutativeSemiringAlg{\mathbb{N}}{+}{\cdot}{0}{1}}{%
    \FormulaRefAuto{NaturalCommutativeUnitalSemiring}[thm]}
  \proofstepwidestar[]{%
    \neg\forall a\in\mathbb{N}\,
    \exists b\in\mathbb{N}\,(a+b=0\land b+a=0)%
  }{\FormulaRefAuto{NaturalNumbersNotRingConclusion}[thm]}
  \proofstepwidestar[]{%
    \CommutativeSemiringAlg{\mathbb{N}}{+}{\cdot}{0}{1}\land
    \neg\forall a\in\mathbb{N}\,
    \exists b\in\mathbb{N}\,(a+b=0\land b+a=0)}{%
    \rAI{1,2}}
\end{tabproofwide}

\end{document}
